%%%%%%%%%%%%%%%%%%%%%%%%%%%%%%%%%%%%%%%%%
% Lachaise Assignment
% LaTeX Template
% Version 1.0 (26/6/2018)
%
% This template originates from:
% http://www.LaTeXTemplates.com
%
% Authors:
% Marion Lachaise & François Févotte
% Vel (vel@LaTeXTemplates.com)
%
% License:
% CC BY-NC-SA 3.0 (http://creativecommons.org/licenses/by-nc-sa/3.0/)
% 
%%%%%%%%%%%%%%%%%%%%%%%%%%%%%%%%%%%%%%%%%

%----------------------------------------------------------------------------------------
%	PACKAGES AND OTHER DOCUMENT CONFIGURATIONS
%----------------------------------------------------------------------------------------

\documentclass[12pt,a4paper]{report}

\input{structure.tex} % Include the file specifying the document structure and custom 

% =========================================================
% Recommended preamble additions for Thai physics notes
% Compile with XeLaTeX
% =========================================================


\usepackage{fontspec}
\usepackage{polyglossia}

\setmainlanguage{thai}
\setotherlanguage{english}

\defaultfontfeatures{Ligatures=TeX}

\setmainfont[
    Path = thsarabun/,
    UprightFont = * ,
    BoldFont = * Bold,
    ItalicFont = * Italic,
    BoldItalicFont = * BoldItalic,
    Extension = .ttf,
    Scale = 1.25
]{THSarabunNew}

\newfontfamily\thaifont[
    Path = thsarabun/,
    UprightFont = * ,
    BoldFont = * Bold,
    ItalicFont = * Italic,
    BoldItalicFont = * BoldItalic,
    Extension = .ttf,
    Scale = 1.25
]{THSarabunNew}

\newfontfamily\thaifontsf[
    Path = thsarabun/,
    UprightFont = * ,
    BoldFont = * Bold,
    ItalicFont = * Italic,
    BoldItalicFont = * BoldItalic,
    Extension = .ttf,
    Scale = 1.25
]{THSarabunNew}

\newfontfamily\thaifonttt[
    Path = thsarabun/,
    UprightFont = * ,
    BoldFont = * Bold,
    ItalicFont = * Italic,
    BoldItalicFont = * BoldItalic,
    Extension = .ttf,
    Scale = 1.25
]{THSarabunNew}

\XeTeXlinebreaklocale "th"
\XeTeXlinebreakskip = 0pt plus 1pt

\usepackage{amsmath, amssymb}
\usepackage{siunitx}
\usepackage{physics}
\usepackage{enumitem}
\usepackage{multicol}
\usepackage{tikz}
\usepackage{pgfplots}
\usepackage[most]{tcolorbox}
\usepackage{graphicx}
\pgfplotsset{compat=1.18}

\sisetup{
    per-mode=symbol,
    detect-all
}

\newtcolorbox{examplebox}[1]{
    colback=blue!3,
    colframe=blue!60!black,
    title=#1,
    fonttitle=\bfseries,
    breakable
}

\newtcolorbox{exercisebox}[1]{
    colback=green!3,
    colframe=green!45!black,
    title=#1,
    fonttitle=\bfseries,
    breakable
}
% =========================================================
% Header and copyright footer
% =========================================================
\usepackage{fancyhdr}

\pagestyle{fancy}
\fancyhf{}

\fancyhead[L]{\small การเคลื่อนที่แนวตรง}
\fancyhead[R]{\small สอวน. ฟิสิกส์ ค่าย 1}

\fancyfoot[L]{\footnotesize \textcopyright\ 2026 Tames Thanakrit Damduan}
\fancyfoot[C]{\footnotesize ใช้เพื่อการเรียนการสอน}
\fancyfoot[R]{\footnotesize หน้า \thepage}

\renewcommand{\headrulewidth}{0.4pt}
\renewcommand{\footrulewidth}{0.4pt}
\begin{document}

% \maketitle % Print the title


% =========================================================
% Chapter 2: Motion in One Dimension
% POSN 1 Physics Preparation Notes
% Language: Thai
% Compiler: XeLaTeX
% =========================================================

\chapter{การเคลื่อนที่แนวตรง (Linear Motion)}
\begin{center}
    {\Large \textbf{Motion in One Dimension}}\\[4pt]
    {\normalsize เอกสารเตรียมสอบคัดเลือก สอวน. ฟิสิกส์ ค่าย 1}\\[6pt]
    {\normalsize จัดทำโดย \textbf{Tames Thanakrit Damduan}}\\[2pt]
    {\footnotesize \textcopyright\ 2026 Tames Thanakrit Damduan. All rights reserved.}\\[8pt]
\end{center}

\section*{เป้าหมายของบทเรียน}

หลังจากเรียนบทนี้ นักเรียนควรทำสิ่งต่อไปนี้ได้

\begin{enumerate}
    \item เข้าใจความหมายของตำแหน่ง การกระจัด ระยะทาง ความเร็ว และความเร่ง
    \item แยกความแตกต่างระหว่างปริมาณสเกลาร์และเวกเตอร์ได้
    \item แยกความแตกต่างระหว่างค่าเฉลี่ยกับค่าขณะหนึ่งได้
    \item ใช้แคลคูลัสพื้นฐานเพื่อหา $v(t)$ และ $a(t)$ จาก $x(t)$ ได้
    \item ใช้กราฟ $x-t$, $v-t$, และ $a-t$ ในการตีความการเคลื่อนที่ได้
    \item ใช้สมการการเคลื่อนที่เมื่อความเร่งคงตัวได้อย่างถูกเงื่อนไข
    \item แก้โจทย์ตกอิสระและโจทย์แนวตรงระดับ สอวน. ได้
\end{enumerate}

\begin{tcolorbox}[title=แนวคิดสำคัญของบทนี้]
การเคลื่อนที่แนวตรงไม่ได้ยากเพราะสูตรเยอะ แต่ยากเพราะนักเรียนมักสับสนเรื่องเครื่องหมาย ทิศทาง และเงื่อนไขการใช้สูตร ดังนั้นทุกข้อให้เริ่มจากกำหนดแกน เลือกทิศบวก เขียนสิ่งที่โจทย์ให้ และเลือกสมการจากความหมายของปริมาณ ไม่ใช่ท่องสูตรอย่างเดียว
\end{tcolorbox}
\newpage{}
% =========================================================
\section{พื้นฐานสำคัญ: สเกลาร์และเวกเตอร์}

ก่อนเรียนเรื่องการเคลื่อนที่ นักเรียนต้องรู้ก่อนว่าปริมาณทางฟิสิกส์บางชนิดบอกแค่ ``มากหรือน้อยเท่าไร'' แต่บางชนิดต้องบอกด้วยว่า ``ไปทางไหน'' ความแตกต่างนี้สำคัญมาก เพราะในบทการเคลื่อนที่แนวตรง เครื่องหมายบวกและลบมักใช้แทนทิศทาง

\subsection{สเกลาร์}

\textbf{แนวคิดก่อนจำชื่อ:} ถ้าปริมาณหนึ่งบอกแค่ขนาดอย่างเดียวก็เพียงพอ ปริมาณนั้นเรียกว่า สเกลาร์

ตัวอย่างเช่น ถ้าบอกว่า

\[
\text{เวลาที่ใช้} = 5\,\si{s}
\]

เราไม่จำเป็นต้องถามว่าเวลาไปทางซ้ายหรือขวา เพราะเวลาไม่มีทิศทาง ดังนั้นเวลาเป็นสเกลาร์

\begin{tcolorbox}[title=นิยาม]
สเกลาร์ (scalar) คือ ปริมาณที่มีเพียงขนาด ไม่มีทิศทาง
\end{tcolorbox}

ตัวอย่างของสเกลาร์ ได้แก่

\[
\text{มวล เวลา อุณหภูมิ พลังงาน ระยะทาง อัตราเร็ว}
\]

\subsection{เวกเตอร์}

\textbf{แนวคิดก่อนจำชื่อ:} ถ้าปริมาณหนึ่งต้องบอกทั้งขนาดและทิศทาง ปริมาณนั้นเรียกว่า เวกเตอร์

ตัวอย่างเช่น ถ้าบอกว่า

\[
\text{รถมีความเร็ว } 20\,\si{m/s}
\]

ยังไม่สมบูรณ์ เพราะยังไม่รู้ว่ารถไปทางไหน แต่ถ้าบอกว่า

\[
\text{รถมีความเร็ว } 20\,\si{m/s} \text{ ไปทางขวา}
\]

ข้อมูลจึงครบทั้งขนาดและทิศทาง

\begin{tcolorbox}[title=นิยาม]
เวกเตอร์ คือ ปริมาณที่มีทั้งขนาดและทิศทาง
\end{tcolorbox}

ตัวอย่างของเวกเตอร์ ได้แก่

\[
\text{การกระจัด ความเร็ว ความเร่ง แรง โมเมนตัม}
\]

\subsection{การแทนทิศทางในแนวตรง}

ในการเคลื่อนที่แนวตรง เราไม่จำเป็นต้องวาดลูกศรทุกครั้ง แต่ใช้เครื่องหมายแทนทิศทางได้ เช่น กำหนดให้

\[
\text{ทิศขวาเป็นบวก}
\]

ดังนั้น

\[
+5\,\si{m}
\]

แปลว่าอยู่ทางขวาหรือเคลื่อนที่ไปทางขวา ส่วน

\[
-5\,\si{m}
\]

แปลว่าอยู่ทางซ้ายหรือเคลื่อนที่ไปทางซ้าย เมื่อเทียบกับแกนที่กำหนด

\begin{center}
\begin{tabular}{|p{0.25\linewidth}|p{0.32\linewidth}|p{0.32\linewidth}|}
\hline
\textbf{หัวข้อ} & \textbf{สเกลาร์} & \textbf{เวกเตอร์} \\[0pt]
\hline
ข้อมูลที่ต้องมี & ขนาดอย่างเดียว & ขนาดและทิศทาง \\
\hline
ตัวอย่าง & ระยะทาง เวลา มวล อัตราเร็ว & การกระจัด ความเร็ว ความเร่ง แรง \\
\hline
เครื่องหมาย & ไม่ได้ใช้แทนทิศทางโดยตรง & เครื่องหมายบวกหรือลบบอกทิศทางได้ในแนวตรง \\
\hline
ตัวอย่างประโยค & รถมีอัตราเร็ว $10\,\si{m/s}$ & รถมีความเร็ว $10\,\si{m/s}$ ไปทางขวา \\
\hline
\end{tabular}
\end{center}

\subsection{ระยะทางกับการกระจัด}

เริ่มจากสถานการณ์ง่าย ๆ นักเรียนเดินไปทางขวา $5\,\si{m}$ แล้วเดินกลับมาทางซ้าย $2\,\si{m}$

สิ่งที่เกิดขึ้นจริงตามเส้นทางคือ เดินทั้งหมด

\[
5+2=7\,\si{m}
\]

ค่านี้เรียกว่า ระยะทาง เพราะสนใจความยาวของเส้นทางทั้งหมด ไม่สนใจทิศทาง

\[
\text{ระยะทาง} = 7\,\si{m}
\]

แต่ถ้าถามว่าตำแหน่งสุดท้ายเปลี่ยนจากตำแหน่งเริ่มต้นไปเท่าไร ให้เทียบจุดเริ่มต้นกับจุดสุดท้ายโดยตรง

\[
\Delta x = 5-2 = 3\,\si{m}
\]

ถ้าเลือกขวาเป็นบวก จะได้ว่า

\[
\Delta x = +3\,\si{m}
\]

ค่านี้เรียกว่า การกระจัด เพราะมีทั้งขนาดและทิศทาง

\begin{tcolorbox}[title=สรุป]
\[
\text{ระยะทาง} = \text{ความยาวของเส้นทางทั้งหมด}
\]

\[
\text{การกระจัด} = \text{ตำแหน่งสุดท้าย} - \text{ตำแหน่งเริ่มต้น}
\]
\end{tcolorbox}

\begin{tcolorbox}[title=ข้อควรจำสำหรับ สอวน.]
คำว่า ``อัตราเร็ว'' เป็นสเกลาร์ แต่คำว่า ``ความเร็ว'' เป็นเวกเตอร์ ดังนั้นความเร็วติดลบได้ แต่อัตราเร็วติดลบไม่ได้
\end{tcolorbox}

\begin{exercisebox}{แบบฝึกหัดสั้น: สเกลาร์และเวกเตอร์}
\begin{enumerate}
    \item จงบอกว่าปริมาณต่อไปนี้เป็นสเกลาร์หรือเวกเตอร์: เวลา, แรง, มวล, ความเร่ง, ระยะทาง, การกระจัด
    \item รถเคลื่อนที่ไปทางขวา $8\,\si{m}$ แล้วกลับซ้าย $3\,\si{m}$ จงหาระยะทางและการกระจัด
    \item ถ้าวัตถุมีความเร็ว $-5\,\si{m/s}$ หมายความว่าอย่างไร
    \item ถ้าวัตถุมีอัตราเร็ว $5\,\si{m/s}$ และอีกวัตถุหนึ่งมีความเร็ว $-5\,\si{m/s}$ สองประโยคนี้ต่างกันอย่างไร
\end{enumerate}
\end{exercisebox}

% =========================================================
\section{คณิตศาสตร์ก่อนเรียน: ลิมิต อนุพันธ์ และพื้นที่ใต้กราฟ}

\subsection{ทำไมต้องใช้แคลคูลัส}

ในการเคลื่อนที่จริง ความเร็วของวัตถุอาจไม่ได้คงที่ เช่น รถออกตัวจากหยุดนิ่งแล้วค่อย ๆ เร็วขึ้น ถ้าเราดูช่วงเวลายาว ๆ เราจะได้แค่ค่าเฉลี่ย แต่ถ้าอยากรู้ว่า ``ตอนนั้นพอดี'' รถเร็วเท่าไร เราต้องใช้แนวคิดของลิมิต

\subsection{ลิมิตจากความเร็วเฉลี่ยไปสู่ความเร็วขณะหนึ่ง}

เริ่มจากสิ่งที่รู้ก่อน คือ ความเร็วเฉลี่ยในช่วงเวลาหนึ่ง

\[
v_{\text{avg}} = \frac{\Delta x}{\Delta t}
\]

ถ้าอยากรู้ความเร็วที่เวลาใดเวลาหนึ่ง เราทำให้ช่วงเวลา $\Delta t$ เล็กลงเรื่อย ๆ

\[
\Delta t \to 0
\]

ดังนั้นความเร็วขณะหนึ่งจึงเป็นค่าที่ความเร็วเฉลี่ยเข้าใกล้เมื่อช่วงเวลาสั้นมาก ๆ

\[
v = \lim_{\Delta t \to 0}\frac{\Delta x}{\Delta t}
\]

ในภาษาของแคลคูลัส ลิมิตนี้คืออนุพันธ์ของตำแหน่งเทียบกับเวลา

\[
v = \frac{dx}{dt}
\]

\begin{tcolorbox}[title=ความหมาย]
อนุพันธ์ในบทนี้ไม่ได้เป็นแค่เครื่องมือคณิตศาสตร์ แต่คือการถามว่า ``ปริมาณหนึ่งเปลี่ยนเร็วแค่ไหนเมื่อเวลาเปลี่ยนไป''
\end{tcolorbox}

\subsection{อนุพันธ์ที่ต้องใช้}

ถ้า $x(t)$ เป็นตำแหน่งของวัตถุ

\[
\text{ตำแหน่ง} \longrightarrow \text{ความเร็ว}
\]

ด้วยการหาอนุพันธ์

\[
v(t)=\frac{dx}{dt}
\]

และถ้า $v(t)$ เป็นความเร็วของวัตถุ

\[
\text{ความเร็ว} \longrightarrow \text{ความเร่ง}
\]

ด้วยการหาอนุพันธ์อีกครั้ง

\[
a(t)=\frac{dv}{dt}
\]

ดังนั้น

\[
a(t)=\frac{d^2x}{dt^2}
\]

\begin{tcolorbox}[title=สูตรอนุพันธ์พื้นฐาน]
\[
\frac{d}{dt}(t^n)=nt^{n-1}
\]

\[
\frac{d}{dt}(c)=0
\]

\[
\frac{d}{dt}(ct^n)=cnt^{n-1}
\]

โดยที่ $c$ เป็นค่าคงตัว
\end{tcolorbox}

\subsection{อินทิกรัลจากพื้นที่ใต้กราฟ}

ตอนหาอนุพันธ์ เราเริ่มจากตำแหน่งแล้วหาอัตราการเปลี่ยนเป็นความเร็ว แต่บางครั้งโจทย์ให้ความเร็วมา แล้วอยากหาการกระจัด เราจึงต้องทำย้อนกลับ

เริ่มจากนิยามของความเร็วในช่วงเวลาสั้นมาก

\[
v \approx \frac{\Delta x}{\Delta t}
\]

จัดรูปเพื่อหาการกระจัดเล็ก ๆ

\[
\Delta x \approx v\Delta t
\]

ถ้าแบ่งเวลาออกเป็นช่วงเล็ก ๆ หลายช่วง การกระจัดรวมคือผลรวมของทุกช่วง

\[
\Delta x \approx \sum v\Delta t
\]

เมื่อทำให้แต่ละช่วงเวลาเล็กมาก ๆ ผลรวมนี้กลายเป็นอินทิกรัล

\[
\Delta x = \int_{t_i}^{t_f} v(t)\,dt
\]

ดังนั้นอินทิกรัลของกราฟ $v-t$ คือพื้นที่ใต้กราฟ และพื้นที่นั้นเท่ากับการกระจัด

\begin{tcolorbox}[title=สรุปความหมายของอินทิกรัล]
\[
\Delta x = \int_{t_i}^{t_f} v(t)\,dt
\]

หมายถึง การกระจัดเท่ากับพื้นที่ใต้กราฟความเร็วกับเวลา
\end{tcolorbox}

ในทำนองเดียวกัน จากนิยามของความเร่ง

\[
a \approx \frac{\Delta v}{\Delta t}
\]

จัดรูปได้

\[
\Delta v \approx a\Delta t
\]

รวมทุกช่วงเวลา

\[
\Delta v = \int_{t_i}^{t_f} a(t)\,dt
\]

\begin{tcolorbox}[title=สูตรอินทิกรัลพื้นฐาน]
\[
\int t^n\,dt=\frac{t^{n+1}}{n+1}+C
\quad (n\neq -1)
\]

\[
\int c\,dt=ct+C
\]
\end{tcolorbox}

\begin{examplebox}{ตัวอย่างคณิตศาสตร์ก่อนเข้าเนื้อหา}
กำหนดให้ตำแหน่งของอนุภาคเป็น

\[
x(t)=3t^2-2t+5
\]

เราต้องการหาความเร็ว จึงใช้ความหมายว่า

\[
v(t)=\frac{dx}{dt}
\]

หาอนุพันธ์ทีละพจน์

\[
\frac{d}{dt}(3t^2)=6t
\]

\[
\frac{d}{dt}(-2t)=-2
\]

\[
\frac{d}{dt}(5)=0
\]

ดังนั้น

\[
v(t)=6t-2
\]

ต่อไปต้องการหาความเร่ง จึงใช้ความหมายว่า

\[
a(t)=\frac{dv}{dt}
\]

\[
a(t)=\frac{d}{dt}(6t-2)=6
\]

ดังนั้นวัตถุมีความเร่งคงตัวเท่ากับ

\[
a=6\,\si{m/s^2}
\]
\end{examplebox}

% =========================================================
\section{ตำแหน่ง ระยะทาง และการกระจัด}

\subsection{ตำแหน่ง}

ก่อนพูดว่าวัตถุเคลื่อนที่อย่างไร เราต้องรู้ก่อนว่าวัตถุอยู่ที่ไหน เมื่อเลือกจุดอ้างอิงและแกนแล้ว เราใช้ $x$ แทนตำแหน่งบนแกนตรง

ตัวอย่างเช่น ถ้าเลือกจุดกำเนิดเป็น $x=0$ และเลือกทิศขวาเป็นบวก

\[
x=+4\,\si{m}
\]

หมายถึงวัตถุอยู่ทางขวาของจุดกำเนิด $4\,\si{m}$ ส่วน

\[
x=-4\,\si{m}
\]

หมายถึงวัตถุอยู่ทางซ้ายของจุดกำเนิด $4\,\si{m}$

\subsection{การกระจัด}

การกระจัดบอกว่าตำแหน่งเปลี่ยนไปเท่าไรจากต้นทางถึงปลายทาง

เริ่มจากนิยามเชิงภาษา

\[
\text{การกระจัด} = \text{ตำแหน่งสุดท้าย} - \text{ตำแหน่งเริ่มต้น}
\]

ใช้สัญลักษณ์ $x_i$ แทนตำแหน่งเริ่มต้น และ $x_f$ แทนตำแหน่งสุดท้าย

\[
\Delta x = x_f-x_i
\]

\begin{examplebox}{ตัวอย่าง}
วัตถุเริ่มที่

\[
x_i=-3\,\si{m}
\]

และไปจบที่

\[
x_f=8\,\si{m}
\]

ดังนั้น

\[
\Delta x=x_f-x_i
\]

แทนค่า

\[
\Delta x=8-(-3)=11\,\si{m}
\]

การกระจัดเป็นบวก แปลว่าตำแหน่งสุดท้ายอยู่ทางบวกเมื่อเทียบกับตำแหน่งเริ่มต้น
\end{examplebox}

\subsection{ระยะทาง}

ระยะทางไม่สนใจทิศทาง แต่สนใจว่าวัตถุเคลื่อนที่จริงตามเส้นทางทั้งหมดกี่เมตร

ถ้าวัตถุไปทางขวา $12\,\si{m}$ แล้วถอยกลับซ้าย $7\,\si{m}$

\[
\text{ระยะทาง}=12+7=19\,\si{m}
\]

แต่การกระจัดคือ

\[
\Delta x=12-7=5\,\si{m}
\]

\begin{tcolorbox}[title=ข้อควรระวัง]
ระยะทางเป็นสเกลาร์ จึงไม่ติดลบ แต่การกระจัดเป็นเวกเตอร์ในแนวตรง จึงเป็นบวก ลบ หรือศูนย์ได้
\end{tcolorbox}

\subsection{แคลคูลัสท้ายหน่วย}

ในหน่วยนี้ให้เริ่มมองตำแหน่งเป็นฟังก์ชันของเวลา เช่น

\[
x(t)=2t^2+3
\]

แปลว่า ตำแหน่งไม่ได้เป็นค่าคงที่ แต่ขึ้นกับเวลา

ถ้าเวลาเปลี่ยนจาก $t=1$ เป็น $t=3$ การกระจัดหาได้จาก

\[
\Delta x=x(3)-x(1)
\]

โดย

\[
x(3)=2(3)^2+3=21
\]

และ

\[
x(1)=2(1)^2+3=5
\]

ดังนั้น

\[
\Delta x=21-5=16\,\si{m}
\]

\begin{exercisebox}{แบบฝึกหัด 2.2}
\begin{enumerate}
    \item อนุภาคเริ่มที่ $x_i=-3\,\si{m}$ และจบที่ $x_f=8\,\si{m}$ จงหาการกระจัด
    \item รถเคลื่อนที่จากตำแหน่ง $x=0$ ไป $x=12\,\si{m}$ แล้วถอยกลับมาที่ $x=5\,\si{m}$ จงหาระยะทางและการกระจัด
    \item วัตถุเคลื่อนที่บนเส้นตรงตามตำแหน่ง $x(t)=t^2-4t+1$ จงหาตำแหน่งเมื่อ $t=0,2,4$
    \item ถ้า $x(t)=3t^2-12t+7$ จงหาช่วงเวลาที่วัตถุกลับมาที่ตำแหน่งเดิมเท่ากับตอน $t=0$
\end{enumerate}
\end{exercisebox}

% =========================================================
\section{ความเร็วเฉลี่ย อัตราเร็วเฉลี่ย และความเร็วขณะหนึ่ง}

\subsection{ความเร็วเฉลี่ย}

เมื่อตำแหน่งของวัตถุเปลี่ยนไป เราต้องการรู้ว่าวัตถุเปลี่ยนตำแหน่งเร็วแค่ไหนในช่วงเวลาหนึ่ง

เริ่มจากความหมาย

\[
\text{ความเร็วเฉลี่ย} = \frac{\text{การกระจัด}}{\text{เวลาที่ใช้}}
\]

เขียนด้วยสัญลักษณ์ได้เป็น

\[
v_{\text{avg}}=\frac{\Delta x}{\Delta t}
\]

และเพราะ

\[
\Delta x=x_f-x_i
\]

\[
\Delta t=t_f-t_i
\]

จึงได้

\[
v_{\text{avg}}=\frac{x_f-x_i}{t_f-t_i}
\]

\begin{examplebox}{ตัวอย่าง}
รถเคลื่อนที่จาก

\[
x_i=2\,\si{m}
\]

ไปยัง

\[
x_f=20\,\si{m}
\]

ในเวลา

\[
\Delta t=6\,\si{s}
\]

เริ่มจากนิยาม

\[
v_{\text{avg}}=\frac{\Delta x}{\Delta t}
\]

หาการกระจัดก่อน

\[
\Delta x=20-2=18\,\si{m}
\]

ดังนั้น

\[
v_{\text{avg}}=\frac{18}{6}=3\,\si{m/s}
\]
\end{examplebox}

\subsection{อัตราเร็วเฉลี่ย}

ถ้าไม่สนใจทิศทาง แต่สนใจว่าวัตถุวิ่งตามเส้นทางจริงเร็วเฉลี่ยเท่าไร ให้ใช้ระยะทางทั้งหมดแทนการกระจัด

เริ่มจากความหมาย

\[
\text{อัตราเร็วเฉลี่ย}=\frac{\text{ระยะทางทั้งหมด}}{\text{เวลาทั้งหมด}}
\]

เขียนเป็นสมการได้ว่า

\[
\text{average speed}=\frac{d}{\Delta t}
\]

โดย $d$ คือระยะทางทั้งหมด

\begin{examplebox}{ตัวอย่าง}
นักเรียนเดินไปทางขวา $10\,\si{m}$ ใน $4\,\si{s}$ แล้วเดินกลับซ้าย $6\,\si{m}$ ใน $2\,\si{s}$

ระยะทางรวมคือ

\[
d=10+6=16\,\si{m}
\]

การกระจัดคือ

\[
\Delta x=10-6=4\,\si{m}
\]

เวลารวมคือ

\[
\Delta t=4+2=6\,\si{s}
\]

อัตราเร็วเฉลี่ยใช้ระยะทาง

\[
\text{average speed}=\frac{16}{6}=\frac{8}{3}\,\si{m/s}
\]

ความเร็วเฉลี่ยใช้การกระจัด

\[
v_{\text{avg}}=\frac{4}{6}=\frac{2}{3}\,\si{m/s}
\]
\end{examplebox}

\subsection{ความเร็วขณะหนึ่ง}

ความเร็วเฉลี่ยบอกภาพรวมของทั้งช่วงเวลา แต่บางครั้งเราต้องการรู้ความเร็ว ณ เวลาหนึ่ง

เริ่มจากนิยามเดิม

\[
v_{\text{avg}}=\frac{\Delta x}{\Delta t}
\]

ถ้าเราทำช่วงเวลาให้เล็กลงเรื่อย ๆ

\[
\Delta t \to 0
\]

ค่าเฉลี่ยในช่วงสั้นมากจะกลายเป็นค่าขณะหนึ่ง

\[
v(t)=\lim_{\Delta t\to 0}\frac{\Delta x}{\Delta t}
\]

ในแคลคูลัส เขียนได้ว่า

\[
v(t)=\frac{dx}{dt}
\]

\begin{tcolorbox}[title=การแปลความหมายของเครื่องหมาย]
ถ้า $v(t)>0$ วัตถุเคลื่อนที่ไปทางบวก

ถ้า $v(t)<0$ วัตถุเคลื่อนที่ไปทางลบ

ถ้า $v(t)=0$ วัตถุหยุดชั่วขณะ แต่อาจไม่ได้หยุดถาวร
\end{tcolorbox}

\begin{examplebox}{ตัวอย่าง}
กำหนดให้

\[
x(t)=t^3-6t^2+9t
\]

ต้องการหาเวลาที่วัตถุหยุดชั่วขณะ เงื่อนไขคือ

\[
v(t)=0
\]

แต่ต้องหา $v(t)$ ก่อนจากนิยาม

\[
v(t)=\frac{dx}{dt}
\]

ดังนั้น

\[
v(t)=3t^2-12t+9
\]

ให้วัตถุหยุดชั่วขณะ

\[
3t^2-12t+9=0
\]

หารด้วย $3$

\[
t^2-4t+3=0
\]

แยกตัวประกอบ

\[
(t-1)(t-3)=0
\]

ดังนั้น

\[
t=1\,\si{s},\quad t=3\,\si{s}
\]
\end{examplebox}

\subsection{แคลคูลัสท้ายหน่วย}

ในกราฟ $x-t$ ความเร็วเฉลี่ยคือความชันของเส้นตรงที่ลากระหว่างสองจุด

\[
v_{\text{avg}}=\frac{x(t_2)-x(t_1)}{t_2-t_1}
\]

แต่ความเร็วขณะหนึ่งคือความชันของเส้นสัมผัสกราฟ ณ เวลานั้น

\[
v(t)=\frac{dx}{dt}
\]

\begin{exercisebox}{แบบฝึกหัด 2.3}
\begin{enumerate}
    \item รถเคลื่อนที่จาก $x=2\,\si{m}$ ไป $x=20\,\si{m}$ ในเวลา $6\,\si{s}$ จงหาความเร็วเฉลี่ย
    \item เด็กเดินทางไปทางขวา $4\,\si{m}$ แล้วเดินกลับซ้าย $4\,\si{m}$ ในเวลา $5\,\si{s}$ จงหาความเร็วเฉลี่ยและอัตราเร็วเฉลี่ย
    \item ถ้า $x(t)=5t^2+2t-1$ จงหา $v(t)$ และ $v(3)$
    \item ถ้า $x(t)=t^3-3t^2-9t$ จงหาเวลาที่วัตถุหยุดชั่วขณะ
    \item วัตถุมีตำแหน่ง $x(t)=2t^3-15t^2+36t$ จงพิจารณาว่าวัตถุเคลื่อนที่ไปทางบวกหรือลบในช่วง $0<t<5$
\end{enumerate}
\end{exercisebox}

% =========================================================
\section{ความเร่ง}

\subsection{ความเร่งเฉลี่ย}

หลังจากรู้ว่าความเร็วคืออัตราการเปลี่ยนตำแหน่ง คำถามต่อมาคือ ความเร็วเปลี่ยนเร็วแค่ไหน ปริมาณนี้เรียกว่า ความเร่ง

เริ่มจากความหมาย

\[
\text{ความเร่งเฉลี่ย}=\frac{\text{การเปลี่ยนความเร็ว}}{\text{เวลาที่ใช้}}
\]

การเปลี่ยนความเร็วคือ

\[
\Delta v=v_f-v_i
\]

ดังนั้น

\[
a_{\text{avg}}=\frac{\Delta v}{\Delta t}
\]

หรือ

\[
a_{\text{avg}}=\frac{v_f-v_i}{t_f-t_i}
\]

\subsection{ความเร่งขณะหนึ่ง}

ถ้าต้องการรู้ความเร่ง ณ เวลาหนึ่ง ให้ทำเหมือนกับตอนเปลี่ยนจากความเร็วเฉลี่ยเป็นความเร็วขณะหนึ่ง

เริ่มจาก

\[
a_{\text{avg}}=\frac{\Delta v}{\Delta t}
\]

ทำให้ช่วงเวลาสั้นลงเรื่อย ๆ

\[
\Delta t\to 0
\]

จึงได้

\[
a(t)=\lim_{\Delta t\to 0}\frac{\Delta v}{\Delta t}
\]

ในแคลคูลัสคือ

\[
a(t)=\frac{dv}{dt}
\]

และเพราะ

\[
v(t)=\frac{dx}{dt}
\]

ดังนั้น

\[
a(t)=\frac{d^2x}{dt^2}
\]

\begin{tcolorbox}[title=กฎจำง่าย]
ถ้า $v$ และ $a$ มีเครื่องหมายเดียวกัน วัตถุเร็วขึ้น

ถ้า $v$ และ $a$ มีเครื่องหมายตรงข้าม วัตถุช้าลง

ถ้า $a=0$ ความเร็วคงตัว

ถ้า $v=0$ ไม่ได้แปลว่า $a=0$
\end{tcolorbox}

\begin{examplebox}{ตัวอย่าง}
วัตถุมีตำแหน่ง

\[
x(t)=2t^3-5t^2+4t
\]

ต้องการหาความเร็ว ใช้

\[
v(t)=\frac{dx}{dt}
\]

ดังนั้น

\[
v(t)=6t^2-10t+4
\]

ต้องการหาความเร่ง ใช้

\[
a(t)=\frac{dv}{dt}
\]

ดังนั้น

\[
a(t)=12t-10
\]
\end{examplebox}

\subsection{แคลคูลัสท้ายหน่วย}

การหาอนุพันธ์เดินหน้าเป็นลำดับนี้

\[
x(t)\xrightarrow{\frac{d}{dt}}v(t)\xrightarrow{\frac{d}{dt}}a(t)
\]

ส่วนการอินทิเกรตเป็นการย้อนกลับ

\[
a(t)\xrightarrow{\int dt}v(t)\xrightarrow{\int dt}x(t)
\]

แต่การอินทิเกรตต้องมีค่าคงตัว ซึ่งต้องหาโดยใช้เงื่อนไขต้น เช่น $x(0)$ หรือ $v(0)$

\begin{examplebox}{ตัวอย่างอินทิเกรต}
กำหนดให้

\[
a(t)=6t
\]

และ

\[
v(0)=4\,\si{m/s}
\]

เริ่มจากความหมาย

\[
a(t)=\frac{dv}{dt}
\]

ดังนั้น

\[
dv=a(t)\,dt
\]

อินทิเกรตทั้งสองข้าง

\[
v(t)=\int 6t\,dt
\]

\[
v(t)=3t^2+C
\]

ใช้เงื่อนไข $v(0)=4$

\[
4=3(0)^2+C
\]

\[
C=4
\]

ดังนั้น

\[
v(t)=3t^2+4
\]
\end{examplebox}

\begin{exercisebox}{แบบฝึกหัด 2.4}
\begin{enumerate}
    \item รถเพิ่มความเร็วจาก $5\,\si{m/s}$ เป็น $25\,\si{m/s}$ ในเวลา $4\,\si{s}$ จงหาความเร่งเฉลี่ย
    \item วัตถุมี $v(t)=4t^2-3t+2$ จงหา $a(t)$
    \item วัตถุมี $x(t)=t^4-4t^2+3$ จงหา $v(t)$ และ $a(t)$
    \item ถ้า $v=-6\,\si{m/s}$ และ $a=2\,\si{m/s^2}$ วัตถุกำลังเร็วขึ้นหรือช้าลง
    \item ถ้า $a(t)=2t+3$ และ $v(0)=1$ จงหา $v(t)$
\end{enumerate}
\end{exercisebox}

% =========================================================
\section{การเคลื่อนที่ด้วยความเร่งคงตัว}

\subsection{เงื่อนไขสำคัญ}

สมการในหน่วยนี้ใช้ได้เมื่อความเร่งคงตัวเท่านั้น

\[
a=\text{constant}
\]

ถ้าความเร่งไม่คงตัว ต้องใช้แคลคูลัสหรือวิธีอื่น ไม่ควรนำสูตรชุดนี้ไปใช้ทันที

\subsection{สมการที่ 1: หาความเร็วจากความเร่งคงตัว}

เริ่มจากนิยามของความเร่ง

\[
a=\frac{\Delta v}{\Delta t}
\]

เมื่อความเร่งคงตัว และให้เริ่มจับเวลาที่ $t=0$

\[
a=\frac{v_f-v_i}{t}
\]

จัดรูป

\[
v_f-v_i=at
\]

ดังนั้น

\[
v_f=v_i+at
\]

\begin{tcolorbox}[title=สมการที่ 1]
\[
v_f=v_i+at
\]
ใช้เมื่อรู้ความเร่งคงตัวและเวลาที่เคลื่อนที่
\end{tcolorbox}

\subsection{สมการที่ 2: หาการกระจัดจากพื้นที่ใต้กราฟ $v-t$}

เมื่อความเร่งคงตัว กราฟ $v-t$ เป็นเส้นตรง เพราะความเร็วเปลี่ยนอย่างสม่ำเสมอ

การกระจัดคือพื้นที่ใต้กราฟ $v-t$

\[
\Delta x=\text{area under }v-t
\]

พื้นที่ใต้กราฟเส้นตรงเป็นรูปสี่เหลี่ยมคางหมู

\[
\Delta x=\frac{1}{2}(v_i+v_f)t
\]

ดังนั้น

\[
x_f-x_i=\frac{v_i+v_f}{2}t
\]

\begin{tcolorbox}[title=สมการที่ 2]
\[
\Delta x=\frac{v_i+v_f}{2}t
\]
ใช้เมื่อรู้ความเร็วต้น ความเร็วปลาย และเวลา
\end{tcolorbox}

\subsection{สมการที่ 3: หาการกระจัดโดยไม่ต้องรู้ความเร็วปลาย}

เริ่มจากสมการที่ 2

\[
\Delta x=\frac{v_i+v_f}{2}t
\]

และใช้สมการที่ 1

\[
v_f=v_i+at
\]

แทนค่า $v_f$ ลงในสมการที่ 2

\[
\Delta x=\frac{v_i+(v_i+at)}{2}t
\]

รวมพจน์ในวงเล็บ

\[
\Delta x=\frac{2v_i+at}{2}t
\]

กระจาย $t$

\[
\Delta x=v_it+\frac{1}{2}at^2
\]

ดังนั้น

\[
x_f=x_i+v_it+\frac{1}{2}at^2
\]

\begin{tcolorbox}[title=สมการที่ 3]
\[
\Delta x=v_it+\frac{1}{2}at^2
\]
ใช้เมื่อรู้ความเร็วต้น ความเร่งคงตัว และเวลา
\end{tcolorbox}

\subsection{สมการที่ 4: หาความเร็วโดยไม่ต้องรู้เวลา}

เริ่มจากสมการที่ 1

\[
v_f=v_i+at
\]

จัดรูปเพื่อหาเวลา

\[
t=\frac{v_f-v_i}{a}
\]

นำไปแทนในสมการที่ 2

\[
\Delta x=\frac{v_i+v_f}{2}t
\]

จะได้

\[
\Delta x=\frac{v_i+v_f}{2}\left(\frac{v_f-v_i}{a}\right)
\]

ใช้ผลต่างกำลังสอง

\[
(v_i+v_f)(v_f-v_i)=v_f^2-v_i^2
\]

ดังนั้น

\[
\Delta x=\frac{v_f^2-v_i^2}{2a}
\]

จัดรูป

\[
v_f^2-v_i^2=2a\Delta x
\]

จึงได้

\[
v_f^2=v_i^2+2a\Delta x
\]

\begin{tcolorbox}[title=สมการที่ 4]
\[
v_f^2=v_i^2+2a\Delta x
\]
ใช้เมื่อไม่ต้องการใช้เวลา
\end{tcolorbox}

\subsection{วิธีเลือกสูตร}

\begin{enumerate}
    \item กำหนดแกนและทิศบวก
    \item เขียนสิ่งที่รู้ เช่น $v_i$, $v_f$, $a$, $t$, $\Delta x$
    \item เขียนสิ่งที่ต้องหา
    \item เลือกสมการที่มีสิ่งที่ต้องหา และไม่มีตัวแปรที่โจทย์ไม่ให้
\end{enumerate}

\begin{examplebox}{ตัวอย่าง}
รถเริ่มจากหยุดนิ่งและเร่งด้วยความเร่งคงตัว $2\,\si{m/s^2}$ เป็นเวลา $5\,\si{s}$ จงหาระยะทาง

เริ่มจากแปลโจทย์

\[
v_i=0,\quad a=2\,\si{m/s^2},\quad t=5\,\si{s}
\]

ต้องการหา

\[
\Delta x
\]

ไม่ต้องรู้ $v_f$ จึงใช้สมการที่ 3

\[
\Delta x=v_it+\frac{1}{2}at^2
\]

แทนค่า

\[
\Delta x=(0)(5)+\frac{1}{2}(2)(5^2)
\]

\[
\Delta x=25\,\si{m}
\]
\end{examplebox}

\subsection{แคลคูลัสท้ายหน่วย: ที่มาของสมการ}

เริ่มจากนิยามของความเร่งขณะหนึ่ง

\[
a=\frac{dv}{dt}
\]

จัดรูป

\[
dv=a\,dt
\]

ถ้า $a$ คงตัว อินทิเกรตจาก $v_i$ ถึง $v_f$ และจาก $0$ ถึง $t$

\[
\int_{v_i}^{v_f}dv=\int_0^t a\,dt
\]

\[
v_f-v_i=at
\]

ดังนั้น

\[
v_f=v_i+at
\]

ต่อไปเริ่มจากนิยามของความเร็ว

\[
v=\frac{dx}{dt}
\]

จัดรูป

\[
dx=v\,dt
\]

เพราะความเร็วในกรณีความเร่งคงตัวคือ

\[
v=v_i+at
\]

แทนลงไป

\[
dx=(v_i+at)\,dt
\]

อินทิเกรต

\[
\int_{x_i}^{x_f}dx=\int_0^t (v_i+at)\,dt
\]

\[
x_f-x_i=v_it+\frac{1}{2}at^2
\]

ดังนั้น

\[
\Delta x=v_it+\frac{1}{2}at^2
\]

\begin{exercisebox}{แบบฝึกหัด 2.5}
\begin{enumerate}
    \item รถเริ่มจากหยุดนิ่ง เร่งด้วย $3\,\si{m/s^2}$ เป็นเวลา $4\,\si{s}$ จงหาความเร็วสุดท้ายและระยะทาง
    \item วัตถุมี $v_i=10\,\si{m/s}$ และชะลอด้วย $a=-2\,\si{m/s^2}$ จงหาเวลาที่หยุด
    \item รถเคลื่อนที่ด้วย $v_i=5\,\si{m/s}$ และมี $a=2\,\si{m/s^2}$ จนความเร็วเป็น $15\,\si{m/s}$ จงหาระยะทาง
    \item วัตถุเคลื่อนที่จากหยุดนิ่งด้วยความเร่งคงตัว ระยะทางในวินาทีที่ 4 เท่ากับเท่าไร ถ้า $a=2\,\si{m/s^2}$
    \item รถ A เริ่มจากหยุดนิ่งด้วยความเร่ง $2a$ รถ B เริ่มพร้อมกันจากหยุดนิ่งด้วยความเร่ง $a$ จงหาอัตราส่วนระยะทางของ A ต่อ B หลังเวลาผ่านไปเท่ากัน
\end{enumerate}
\end{exercisebox}

% =========================================================
\section{การตกอิสระ}

\subsection{แนวคิดหลัก}

การตกอิสระคือการเคลื่อนที่ที่มีแรงโน้มถ่วงเป็นสาเหตุหลักของความเร่ง โดยละเลยแรงต้านอากาศ ใกล้ผิวโลกใช้ค่าประมาณ

\[
g\approx 9.8\,\si{m/s^2}
\]

ทิศของความเร่งโน้มถ่วงชี้ลงเสมอ

ถ้าเลือกแกนขึ้นเป็นบวก ความเร่งคือ

\[
a=-g
\]

ถ้าเลือกแกนลงเป็นบวก ความเร่งคือ

\[
a=+g
\]

\begin{tcolorbox}[title=จุดสำคัญ]
ที่จุดสูงสุดของการโยนขึ้นในแนวดิ่ง ความเร็วเป็นศูนย์ชั่วขณะ แต่ความเร่งยังชี้ลง ดังนั้นถ้าเลือกขึ้นเป็นบวก ความเร่งยังเป็น $-g$
\end{tcolorbox}

\subsection{สมการตกอิสระจากสมการความเร่งคงตัว}

การตกอิสระใกล้ผิวโลกเป็นกรณีของความเร่งคงตัว ดังนั้นเริ่มจากสมการเดิม

\[
v_f=v_i+at
\]

ถ้าเลือกแกนขึ้นเป็นบวก

\[
a=-g
\]

แทนค่าได้

\[
v_f=v_i-gt
\]

เริ่มจากสมการตำแหน่ง

\[
\Delta y=v_it+\frac{1}{2}at^2
\]

แทน $a=-g$

\[
\Delta y=v_it-\frac{1}{2}gt^2
\]

หรือ

\[
y_f=y_i+v_it-\frac{1}{2}gt^2
\]

\begin{examplebox}{ตัวอย่าง}
โยนวัตถุขึ้นด้วยความเร็วต้น $20\,\si{m/s}$ จงหาความสูงสูงสุด ใช้ $g=10\,\si{m/s^2}$

เลือกแกนขึ้นเป็นบวก

\[
v_i=20\,\si{m/s}
\]

ที่จุดสูงสุด วัตถุหยุดชั่วขณะ

\[
v_f=0
\]

ความเร่งคือ

\[
a=-10\,\si{m/s^2}
\]

ต้องการหาความสูง จึงใช้สมการที่ไม่มีเวลา

\[
v_f^2=v_i^2+2a\Delta y
\]

แทนค่า

\[
0^2=20^2+2(-10)\Delta y
\]

\[
0=400-20\Delta y
\]

\[
20\Delta y=400
\]

\[
\Delta y=20\,\si{m}
\]
\end{examplebox}

\subsection{แคลคูลัสท้ายหน่วย}

ถ้าเลือกแกนขึ้นเป็นบวก ความเร่งเป็นค่าคงตัว

\[
a(t)=-g
\]

เริ่มจาก

\[
a=\frac{dv}{dt}
\]

จึงได้

\[
dv=-g\,dt
\]

อินทิเกรต

\[
v(t)=\int -g\,dt
\]

\[
v(t)=-gt+C
\]

ใช้เงื่อนไข $v(0)=v_0$

\[
v_0=C
\]

ดังนั้น

\[
v(t)=v_0-gt
\]

ต่อไปเริ่มจาก

\[
v=\frac{dy}{dt}
\]

จึงได้

\[
dy=v\,dt
\]

แทน $v(t)=v_0-gt$

\[
dy=(v_0-gt)\,dt
\]

อินทิเกรต

\[
y(t)=v_0t-\frac{1}{2}gt^2+C_2
\]

ใช้เงื่อนไข $y(0)=y_0$

\[
C_2=y_0
\]

ดังนั้น

\[
y(t)=y_0+v_0t-\frac{1}{2}gt^2
\]

\begin{exercisebox}{แบบฝึกหัด 2.6}
\begin{enumerate}
    \item ปล่อยวัตถุจากหยุดนิ่งจากความสูง $45\,\si{m}$ ใช้ $g=10\,\si{m/s^2}$ จงหาเวลาตกถึงพื้น
    \item โยนวัตถุขึ้นด้วย $30\,\si{m/s}$ ใช้ $g=10\,\si{m/s^2}$ จงหาความสูงสูงสุด
    \item จากข้อก่อนหน้า จงหาเวลารวมที่วัตถุกลับมาระดับเดิม
    \item ปล่อยหินจากตึกสูง $80\,\si{m}$ และหลังจากนั้น $1\,\si{s}$ โยนหินอีกก้อนลงด้วยความเร็วต้น $10\,\si{m/s}$ จงหาว่าก้อนใดถึงพื้นก่อน
    \item วัตถุถูกโยนขึ้นจากพื้น เมื่อเวลาผ่านไป $2\,\si{s}$ อยู่สูงจากพื้น $30\,\si{m}$ และเมื่อเวลาผ่านไป $4\,\si{s}$ อยู่สูงจากพื้น $40\,\si{m}$ ใช้ $g=10\,\si{m/s^2}$ จงหาความเร็วต้น
\end{enumerate}
\end{exercisebox}

% =========================================================
\section{กราฟการเคลื่อนที่}

\subsection{กราฟตำแหน่งกับเวลา}

เริ่มจากนิยามของความเร็วเฉลี่ย

\[
v_{\text{avg}}=\frac{\Delta x}{\Delta t}
\]

ในกราฟ $x-t$ ค่า

\[
\frac{\Delta x}{\Delta t}
\]

คือความชันของกราฟ ดังนั้นความชันของกราฟ $x-t$ คือความเร็วเฉลี่ยในช่วงนั้น

ถ้าดูช่วงเวลาสั้นมาก ๆ ความชันของเส้นสัมผัสคือความเร็วขณะหนึ่ง

\[
v=\frac{dx}{dt}
\]

\begin{tcolorbox}[title=สรุปกราฟ $x-t$]
ความชันของกราฟ $x-t$ คือความเร็ว
\end{tcolorbox}

\subsection{กราฟความเร็วกับเวลา}

เริ่มจากนิยามของความเร่งเฉลี่ย

\[
a_{\text{avg}}=\frac{\Delta v}{\Delta t}
\]

ในกราฟ $v-t$ ค่า

\[
\frac{\Delta v}{\Delta t}
\]

คือความชันของกราฟ ดังนั้นความชันของกราฟ $v-t$ คือความเร่ง

\[
a=\frac{dv}{dt}
\]

ต่อไปเริ่มจากความหมายของความเร็ว

\[
v\approx \frac{\Delta x}{\Delta t}
\]

จัดรูป

\[
\Delta x\approx v\Delta t
\]

บนกราฟ $v-t$ ค่านี้คือพื้นที่ของแถบเล็ก ๆ ใต้กราฟ เมื่อรวมทุกแถบจึงได้

\[
\Delta x=\int v\,dt
\]

\begin{tcolorbox}[title=สรุปกราฟ $v-t$]
ความชันของกราฟ $v-t$ คือความเร่ง

พื้นที่ใต้กราฟ $v-t$ คือการกระจัด
\end{tcolorbox}

\subsection{กราฟความเร่งกับเวลา}

เริ่มจากความหมายของความเร่ง

\[
a\approx \frac{\Delta v}{\Delta t}
\]

จัดรูป

\[
\Delta v\approx a\Delta t
\]

บนกราฟ $a-t$ ค่านี้คือพื้นที่ของแถบเล็ก ๆ ใต้กราฟ เมื่อรวมทุกแถบจึงได้

\[
\Delta v=\int a\,dt
\]

\begin{tcolorbox}[title=สรุปกราฟ $a-t$]
พื้นที่ใต้กราฟ $a-t$ คือการเปลี่ยนความเร็ว
\end{tcolorbox}

\subsection{แคลคูลัสท้ายหน่วย}

ความสัมพันธ์ของกราฟและแคลคูลัสสรุปได้ดังนี้

\[
v=\frac{dx}{dt}
\]

แปลว่า $v$ คือความชันของกราฟ $x-t$

\[
a=\frac{dv}{dt}
\]

แปลว่า $a$ คือความชันของกราฟ $v-t$

\[
\Delta x=\int v\,dt
\]

แปลว่า $\Delta x$ คือพื้นที่ใต้กราฟ $v-t$

\[
\Delta v=\int a\,dt
\]

แปลว่า $\Delta v$ คือพื้นที่ใต้กราฟ $a-t$
\begin{exercisebox}{แบบฝึกหัด 2.7}
\begin{enumerate}
    \item กราฟ $v-t$ เป็นเส้นตรงจาก $v=0$ ถึง $v=20\,\si{m/s}$ ใน $5\,\si{s}$ จงหาการกระจัด
    \item กราฟ $a-t$ มี $a=4\,\si{m/s^2}$ คงตัวเป็นเวลา $3\,\si{s}$ ถ้า $v_i=2\,\si{m/s}$ จงหา $v_f$
    \item ถ้า $v(t)=6t-4$ จงหาการกระจัดจาก $t=1$ ถึง $t=3$
    \item ถ้า $a(t)=2t$ และ $v(0)=3$ จงหา $v(4)$
    \item ถ้า $x(t)=t^3-6t^2+12t$ จงวิเคราะห์ช่วงที่วัตถุเร็วขึ้นและช้าลงในช่วง $0<t<5$
\end{enumerate}
\end{exercisebox}

% =========================================================
\section{ชุดโจทย์รวมท้ายบท: คละระดับสำหรับ สอวน. ค่าย 1}

\subsection*{ระดับพื้นฐาน}

\begin{enumerate}
    \item วัตถุเคลื่อนที่จาก $x=-5\,\si{m}$ ไป $x=15\,\si{m}$ ใน $4\,\si{s}$ จงหาความเร็วเฉลี่ย
    \item วัตถุเริ่มจากหยุดนิ่งด้วยความเร่ง $2\,\si{m/s^2}$ เป็นเวลา $10\,\si{s}$ จงหาระยะทาง
    \item โยนวัตถุขึ้นด้วย $20\,\si{m/s}$ ใช้ $g=10\,\si{m/s^2}$ จงหาเวลาถึงจุดสูงสุด
    \item วัตถุมี $x(t)=4t^2-3t+1$ จงหา $v(t)$ และ $a(t)$
\end{enumerate}

\subsection*{ระดับกลาง}

\begin{enumerate}
    \setcounter{enumi}{4}
    \item รถเริ่มจากหยุดนิ่งและวิ่งด้วยความเร่งคงตัว ระยะทางในวินาทีที่ 5 เท่ากับ $18\,\si{m}$ จงหาความเร่ง
    \item วัตถุมี $v(t)=t^2-4t+3$ จงหาการกระจัดจาก $t=0$ ถึง $t=5$
    \item โยนวัตถุขึ้นจากพื้นด้วยความเร็วต้น $u$ จงพิสูจน์ว่าความสูงสูงสุดคือ
    \[
    H=\frac{u^2}{2g}
    \]
    \item รถสองคันเริ่มจากจุดเดียวกัน รถ A เริ่มจากหยุดนิ่งด้วยความเร่ง $a$ รถ B ออกตัวด้วยความเร็วคงตัว $u$ จงหาเวลาที่ A ตามทัน B
\end{enumerate}

\subsection*{ระดับยาก}

\begin{enumerate}
    \setcounter{enumi}{8}
    \item วัตถุเคลื่อนที่ตาม
    \[
    x(t)=t^3-9t^2+24t
    \]
    จงหาเวลาที่วัตถุเปลี่ยนทิศทาง และหาตำแหน่งที่เกิดเหตุการณ์นั้น
    \item วัตถุเริ่มที่ $x=0$ และ $v=0$ โดยมีความเร่ง
    \[
    a(t)=6t-4
    \]
    จงหา $x(t)$
    \item ลิฟต์กำลังเคลื่อนที่ขึ้นด้วยความเร็ว $v_0$ เมื่อวัตถุหลุดจากมือคนในลิฟต์ จงเขียนสมการตำแหน่งของวัตถุเทียบกับพื้นโลก โดยเลือกแกนขึ้นเป็นบวก
    \item ลูกบอลถูกโยนขึ้นจากพื้นด้วยความเร็วต้น $u$ หลังจากเวลาผ่านไป $T$ ลูกบอลอยู่ที่ความสูง $h$ จงหาค่า $u$ ในรูปของ $h,T,g$
    \item รถคันหนึ่งเคลื่อนที่ด้วยความเร็วต้น $u$ แล้วเบรกด้วยความเร่งขนาด $a$ จนหยุด ระยะหยุดคือ $s$ ถ้าความเร็วต้นเพิ่มเป็น $2u$ และความเร่งเบรกมีขนาดเท่าเดิม ระยะหยุดเป็นกี่เท่าของเดิม
    \item วัตถุเคลื่อนที่โดยมี
    \[
    v(x)=kx
    \]
    เมื่อ $k$ เป็นค่าคงตัวบวก จงหา $x(t)$ ถ้า $x(0)=x_0$
\end{enumerate}

% =========================================================
\section{เฉลยย่อท้ายบท}

\begin{enumerate}
    \item $v_{\text{avg}}=5\,\si{m/s}$
    \item $\Delta x=100\,\si{m}$
    \item $t=2\,\si{s}$
    \item $v(t)=8t-3,\quad a=8\,\si{m/s^2}$
    \item ระยะทางในวินาทีที่ $n$ จากหยุดนิ่งคือ
    \[
    s_n=\frac{1}{2}a(2n-1)
    \]
    ดังนั้น
    \[
    18=\frac{1}{2}a(9)
    \Rightarrow a=4\,\si{m/s^2}
    \]
    \item
    \[
    \Delta x=\int_0^5 (t^2-4t+3)\,dt
    =\left[\frac{t^3}{3}-2t^2+3t\right]_0^5
    =\frac{125}{3}-50+15
    =\frac{20}{3}\,\si{m}
    \]
    \item ที่จุดสูงสุด $v=0$
    \[
    0=u^2-2gH
    \Rightarrow H=\frac{u^2}{2g}
    \]
    \item
    \[
    x_A=\frac{1}{2}at^2,\quad x_B=ut
    \]
    ตามทันเมื่อ
    \[
    \frac{1}{2}at^2=ut
    \Rightarrow t=\frac{2u}{a}
    \]
    โดยไม่นับ $t=0$
    \item
    \[
    v(t)=3t^2-18t+24=3(t-2)(t-4)
    \]
    เปลี่ยนทิศที่ $t=2,4$
    \[
    x(2)=8-36+48=20
    \]
    \[
    x(4)=64-144+96=16
    \]
    \item
    \[
    v(t)=\int(6t-4)\,dt=3t^2-4t+C
    \]
    เพราะ $v(0)=0$ จึง $C=0$
    \[
    x(t)=\int(3t^2-4t)\,dt=t^3-2t^2+C_2
    \]
    เพราะ $x(0)=0$ จึง
    \[
    x(t)=t^3-2t^2
    \]
    \item
    \[
    y(t)=y_0+v_0t-\frac{1}{2}gt^2
    \]
    โดย $v_0$ เป็นความเร็วของลิฟต์ขณะปล่อย
    \item
    \[
    h=uT-\frac{1}{2}gT^2
    \Rightarrow
    u=\frac{h}{T}+\frac{1}{2}gT
    \]
    \item
    \[
    s=\frac{u^2}{2a}
    \]
    ถ้าเปลี่ยนเป็น $2u$
    \[
    s'=\frac{(2u)^2}{2a}=4s
    \]
    ดังนั้นเป็น $4$ เท่า
    \item
    จาก
    \[
    v=\frac{dx}{dt}=kx
    \]
    แยกตัวแปร
    \[
    \frac{dx}{x}=k\,dt
    \]
    อินทิเกรต
    \[
    \ln x=kt+C
    \]
    ใช้ $x(0)=x_0$
    \[
    x(t)=x_0e^{kt}
    \]
\end{enumerate}

% =========================================================
\section{สรุปสูตรสำคัญ}

\begin{tcolorbox}[title=สรุปบทที่ 1]
\[
\Delta x=x_f-x_i
\]

\[
v_{\text{avg}}=\frac{\Delta x}{\Delta t}
\]

\[
v=\frac{dx}{dt}
\]

\[
a_{\text{avg}}=\frac{\Delta v}{\Delta t}
\]

\[
a=\frac{dv}{dt}=\frac{d^2x}{dt^2}
\]

เมื่อ $a$ คงตัว

\[
v_f=v_i+at
\]

\[
\Delta x=v_it+\frac{1}{2}at^2
\]

\[
v_f^2=v_i^2+2a\Delta x
\]

\[
\Delta x=\frac{v_i+v_f}{2}t
\]

ตกอิสระใกล้ผิวโลก

\[
a=-g
\quad \text{เมื่อเลือกขึ้นเป็นบวก}
\]
\end{tcolorbox}


% =========================================================
% Chapter 3: Vectors
% POSN 1 Physics Preparation Notes
% Language: Thai
% Compiler: XeLaTeX
% =========================================================

\chapter{เวกเตอร์}

\begin{center}
    {\Large \textbf{Vectors}}\\[4pt]
    {\normalsize เอกสารเตรียมสอบคัดเลือก สอวน. ฟิสิกส์ ค่าย 1}\\[6pt]
    {\normalsize จัดทำโดย \textbf{Tames Thanakrit Damduan}}\\[2pt]
    {\footnotesize \textcopyright\ 2026 Tames Thanakrit Damduan. All rights reserved.}\\[8pt]
\end{center}

\section*{เป้าหมายของบทเรียน}

หลังจากเรียนบทนี้ นักเรียนควรทำสิ่งต่อไปนี้ได้

\begin{enumerate}
    \item เข้าใจความแตกต่างระหว่างปริมาณสเกลาร์และเวกเตอร์
    \item ใช้ระบบพิกัดฉากและระบบพิกัดเชิงขั้วในการระบุตำแหน่งได้
    \item แปลงพิกัดระหว่างรูป $(x,y)$ และรูป $(r,\theta)$ ได้
    \item บวก ลบ และคูณเวกเตอร์ด้วยจำนวนจริงได้
    \item แตกเวกเตอร์เป็นองค์ประกอบตามแกน $x$ และ $y$ ได้
    \item เขียนเวกเตอร์ในรูปเวกเตอร์หน่วย $\hat{i}$, $\hat{j}$, และ $\hat{k}$ ได้
    \item ใช้เวกเตอร์ในการเตรียมต่อยอดไปยังแรง โมเมนตัม สนามไฟฟ้า และโพรเจคไทล์ได้
\end{enumerate}

\begin{tcolorbox}[title=แนวคิดสำคัญของบทนี้]
เวกเตอร์ไม่ใช่แค่ลูกศร แต่คือภาษาสำหรับบอกทั้งขนาดและทิศทางของปริมาณทางฟิสิกส์ จุดที่สำคัญที่สุดคือการแตกเวกเตอร์เป็นองค์ประกอบตามแกน เพราะหลังจากนั้นโจทย์สองมิติจะกลายเป็นโจทย์หนึ่งมิติหลายแกน
\end{tcolorbox}

% =========================================================
\section{คณิตศาสตร์ก่อนเรียน: พิกัด มุม และตรีโกณมิติ}

\subsection{ทำไมต้องมีระบบพิกัด}

ในบทที่แล้ว เราใช้เพียงแกนเดียว เช่น แกน $x$ เพื่อบอกตำแหน่งของวัตถุ แต่เมื่อวัตถุอยู่ในระนาบ เช่น บนกระดาน บนพื้น หรือในอากาศ เราต้องบอกตำแหน่งด้วยข้อมูลมากกว่าหนึ่งค่า

ตัวอย่างเช่น จุดหนึ่งบนระนาบอาจอยู่ทางขวา $3\,\si{m}$ และอยู่สูงขึ้น $4\,\si{m}$ จากจุดกำเนิด เราจึงเขียนเป็น

\[
(x,y)=(3,4)
\]

ระบบนี้เรียกว่า ระบบพิกัดฉาก หรือ Cartesian coordinate system

\begin{tcolorbox}[title=นิยาม]
ระบบพิกัดฉาก คือ ระบบที่ใช้แกนตั้งฉากกัน เช่น แกน $x$ และแกน $y$ เพื่อบอกตำแหน่งของจุดในระนาบ
\end{tcolorbox}

\subsection{พิกัดฉาก}

ถ้าจุด $P$ มีพิกัด

\[
P=(x,y)
\]

หมายความว่า จากจุดกำเนิด $O$ ต้องเลื่อนไปตามแกน $x$ เป็นระยะ $x$ แล้วเลื่อนไปตามแกน $y$ เป็นระยะ $y$

ตัวอย่างเช่น

\[
P=(5,3)
\]

แปลว่า จุดนี้อยู่ทางขวา $5$ หน่วย และอยู่สูงขึ้น $3$ หน่วย

\[
Q=(-3,4)
\]

แปลว่า จุดนี้อยู่ทางซ้าย $3$ หน่วย และอยู่สูงขึ้น $4$ หน่วย

\subsection{พิกัดเชิงขั้ว}

บางครั้งการบอกตำแหน่งด้วยระยะทางจากจุดกำเนิดและมุมจะสะดวกกว่า เช่น ต้องบอกทิศทางของเรือ เครื่องบิน หรือแรง

เราใช้พิกัดเชิงขั้วในรูป

\[
(r,\theta)
\]

โดยที่

\[
r = \text{ระยะจากจุดกำเนิดถึงจุดนั้น}
\]

\[
\theta = \text{มุมที่วัดจากแกน } +x \text{ ทวนเข็มนาฬิกา}
\]

\begin{tcolorbox}[title=ข้อควรจำ]
ในเอกสารนี้ ถ้าไม่ได้บอกเป็นอย่างอื่น จะถือว่า $\theta$ วัดจากแกน $+x$ ทวนเข็มนาฬิกา
\end{tcolorbox}

\subsection{แปลงจากพิกัดเชิงขั้วเป็นพิกัดฉาก}

เริ่มจากรูปสามเหลี่ยมมุมฉากที่มีด้านตรงข้ามมุมฉากยาว $r$ และทำมุม $\theta$ กับแกน $x$

จากนิยามของโคไซน์

\[
\cos\theta=\frac{\text{ด้านประชิดมุม}}{\text{ด้านตรงข้ามมุมฉาก}}
\]

ในที่นี้ด้านประชิดมุมคือ $x$ และด้านตรงข้ามมุมฉากคือ $r$

\[
\cos\theta=\frac{x}{r}
\]

จัดรูปได้

\[
x=r\cos\theta
\]

จากนิยามของไซน์

\[
\sin\theta=\frac{\text{ด้านตรงข้ามมุม}}{\text{ด้านตรงข้ามมุมฉาก}}
\]

ในที่นี้ด้านตรงข้ามมุมคือ $y$

\[
\sin\theta=\frac{y}{r}
\]

จัดรูปได้

\[
y=r\sin\theta
\]

\begin{tcolorbox}[title={สูตรแปลงจาก $(r,\theta)$ ไปเป็น $(x,y)$}]
\[
x=r\cos\theta
\]

\[
y=r\sin\theta
\]
\end{tcolorbox}

\subsection{แปลงจากพิกัดฉากเป็นพิกัดเชิงขั้ว}

เริ่มจากจุดที่มีพิกัด $(x,y)$ เราต้องการหาระยะ $r$ จากจุดกำเนิดถึงจุดนั้น

จากทฤษฎีพีทาโกรัส

\[
r^2=x^2+y^2
\]

ดังนั้น

\[
r=\sqrt{x^2+y^2}
\]

ต่อไปต้องการหามุม $\theta$

จากนิยามของแทนเจนต์

\[
\tan\theta=\frac{\text{ด้านตรงข้ามมุม}}{\text{ด้านประชิดมุม}}
\]

ในที่นี้ด้านตรงข้ามคือ $y$ และด้านประชิดคือ $x$

\[
\tan\theta=\frac{y}{x}
\]

ดังนั้น

\[
\theta=\tan^{-1}\left(\frac{y}{x}\right)
\]

\begin{tcolorbox}[title={สูตรแปลงจาก $(x,y)$ ไปเป็น $(r,\theta)$}]
\[
r=\sqrt{x^2+y^2}
\]

\[
\tan\theta=\frac{y}{x}
\]
\end{tcolorbox}

\begin{tcolorbox}[title=จุดที่นักเรียนมักพลาด]
สูตร $\theta=\tan^{-1}(y/x)$ ต้องระวังควอดแรนต์ เพราะเครื่องคิดเลขอาจให้มุมหลักที่ไม่ได้บอกทิศทางจริงของเวกเตอร์ ต้องดูเครื่องหมายของ $x$ และ $y$ ประกอบเสมอ
\end{tcolorbox}

\begin{examplebox}{ตัวอย่าง: แปลงพิกัดฉากเป็นพิกัดเชิงขั้ว}
จุด $P$ มีพิกัด

\[
P=(3,4)
\]

ต้องการหา $r$ เริ่มจาก

\[
r=\sqrt{x^2+y^2}
\]

แทนค่า

\[
r=\sqrt{3^2+4^2}
\]

\[
r=5
\]

ต้องการหามุม

\[
\tan\theta=\frac{y}{x}
\]

\[
\tan\theta=\frac{4}{3}
\]

ดังนั้น

\[
\theta=\tan^{-1}\left(\frac{4}{3}\right)
\]

จุดนี้อยู่ควอดแรนต์ที่ 1 เพราะ $x>0$ และ $y>0$
\end{examplebox}

\begin{exercisebox}{แบบฝึกหัดคณิตศาสตร์ก่อนเรียน}
\begin{enumerate}
    \item จงแปลงจุด $(6,8)$ ให้อยู่ในรูป $(r,\theta)$
    \item จงแปลงจุด $(r,\theta)=(10,30^\circ)$ ให้อยู่ในรูป $(x,y)$
    \item จุดหนึ่งมี $x=-3$ และ $y=4$ จงหา $r$ และระบุว่าจุดนี้อยู่ควอดแรนต์ใด
    \item เวกเตอร์หนึ่งยาว $20\,\si{N}$ ทำมุม $60^\circ$ กับแกน $+x$ จงหาขนาดองค์ประกอบตามแกน $x$ และ $y$
\end{enumerate}
\end{exercisebox}

% =========================================================
\section{หน่วยที่ 1: สเกลาร์และเวกเตอร์}

\subsection{เริ่มจากคำถามง่าย ๆ}

ถ้าบอกว่า

\[
\text{ระยะทางจากโรงเรียนถึงบ้านคือ } 5\,\si{km}
\]

ข้อมูลนี้มีแค่ขนาด และเพียงพอสำหรับคำว่า ``ระยะทาง''

แต่ถ้าบอกว่า

\[
\text{โรงเรียนอยู่ห่างจากบ้าน } 5\,\si{km}
\]

ยังไม่พอสำหรับการเดินทางจริง เพราะต้องรู้ด้วยว่าอยู่ทางไหน เช่น ทิศเหนือ ทิศตะวันออก หรือทำมุมเท่าไรกับถนน

นี่คือเหตุผลที่ฟิสิกส์ต้องแยกปริมาณออกเป็นสองชนิด คือ สเกลาร์และเวกเตอร์

\subsection{สเกลาร์}

\begin{tcolorbox}[title=นิยาม]
สเกลาร์ คือ ปริมาณที่มีเพียงขนาด ไม่มีทิศทาง
\end{tcolorbox}

ตัวอย่างเช่น

\[
\text{เวลา มวล อุณหภูมิ พลังงาน ระยะทาง อัตราเร็ว}
\]

สเกลาร์สามารถบวก ลบ คูณ หารเหมือนจำนวนทั่วไปได้ แต่ต้องเป็นปริมาณชนิดเดียวกันและมีหน่วยเข้ากันได้

ตัวอย่างเช่น

\[
3\,\si{s}+5\,\si{s}=8\,\si{s}
\]

แต่

\[
3\,\si{s}+5\,\si{m}
\]

ไม่มีความหมายทางฟิสิกส์ เพราะเป็นคนละชนิดของปริมาณ

\subsection{เวกเตอร์}

\begin{tcolorbox}[title=นิยาม]
เวกเตอร์ คือ ปริมาณที่มีทั้งขนาดและทิศทาง
\end{tcolorbox}

ตัวอย่างเช่น

\[
\text{การกระจัด ความเร็ว ความเร่ง แรง โมเมนตัม สนามไฟฟ้า}
\]

เวกเตอร์มักเขียนด้วยสัญลักษณ์มีลูกศร เช่น

\[
\vec{A}
\]

หรือในหนังสือบางเล่มอาจใช้ตัวหนา เช่น $\mathbf{A}$

ขนาดของเวกเตอร์เขียนเป็น

\[
|\vec{A}|=A
\]

ซึ่งเป็นสเกลาร์และไม่ติดลบ

\subsection{เวกเตอร์ที่เท่ากัน}

เวกเตอร์สองตัวจะเท่ากันเมื่อมี

\[
\text{ขนาดเท่ากัน}
\]

และ

\[
\text{ทิศทางเดียวกัน}
\]

ตำแหน่งที่วาดไม่สำคัญ เพราะเวกเตอร์สามารถเลื่อนขนานไปวางที่อื่นได้โดยยังเป็นเวกเตอร์เดิม ตราบใดที่ขนาดและทิศทางไม่เปลี่ยน

\begin{tcolorbox}[title=ข้อควรจำ]
เวกเตอร์ไม่จำเป็นต้องเริ่มที่จุดกำเนิดเสมอไป สิ่งที่นิยามเวกเตอร์คือขนาดและทิศทาง ไม่ใช่ตำแหน่งที่วาด
\end{tcolorbox}

\subsection{เวกเตอร์ตรงข้าม}

เวกเตอร์ตรงข้ามของ $\vec{A}$ เขียนเป็น

\[
-\vec{A}
\]

มีขนาดเท่ากับ $\vec{A}$ แต่มีทิศตรงข้าม

ดังนั้น

\[
\vec{A}+(-\vec{A})=\vec{0}
\]

\begin{exercisebox}{แบบฝึกหัดหน่วยที่ 1}
\begin{enumerate}
    \item จงจำแนกปริมาณต่อไปนี้ว่าเป็นสเกลาร์หรือเวกเตอร์: แรง, พลังงาน, ความเร็ว, อัตราเร็ว, สนามไฟฟ้า, เวลา
    \item เวกเตอร์สองตัวมีขนาดเท่ากัน แต่ตัวหนึ่งชี้ทิศตะวันออก อีกตัวหนึ่งชี้ทิศตะวันตก เวกเตอร์สองตัวนี้เท่ากันหรือไม่
    \item ถ้า $\vec{A}$ มีขนาด $12\,\si{N}$ ไปทางขวา จงอธิบายความหมายของ $-\vec{A}$
    \item จงยกตัวอย่างสถานการณ์จริงที่ต้องใช้เวกเตอร์ในการอธิบาย ไม่สามารถใช้สเกลาร์อย่างเดียวได้
\end{enumerate}
\end{exercisebox}

% =========================================================
\section{หน่วยที่ 2: การบวกและลบเวกเตอร์}

\subsection{เริ่มจากสถานการณ์การเดินทาง}

นักเรียนเดินจากจุดเริ่มต้นไปทางตะวันออก $3\,\si{m}$ แล้วเดินไปทางเหนือ $4\,\si{m}$

ถ้าถามว่าเดินเป็นระยะทางรวมเท่าไร คำตอบคือ

\[
3+4=7\,\si{m}
\]

แต่ถ้าถามว่าตำแหน่งสุดท้ายห่างจากจุดเริ่มต้นเท่าไร คำตอบไม่ใช่ $7\,\si{m}$ เพราะทิศทางทั้งสองตั้งฉากกัน

การกระจัดลัพธ์คือเส้นตรงจากจุดเริ่มต้นถึงจุดสุดท้าย

\[
R=\sqrt{3^2+4^2}=5\,\si{m}
\]

นี่แสดงว่า การบวกเวกเตอร์ไม่เหมือนการบวกจำนวนธรรมดา

\subsection{วิธีหัวต่อหาง}

ให้ต้องการบวก

\[
\vec{R}=\vec{A}+\vec{B}
\]

วิธีคิดคือ

\begin{enumerate}
    \item วาดเวกเตอร์ $\vec{A}$
    \item นำหางของเวกเตอร์ $\vec{B}$ ไปต่อที่หัวของ $\vec{A}$
    \item เวกเตอร์ลัพธ์ $\vec{R}$ คือเวกเตอร์จากหางของ $\vec{A}$ ไปยังหัวของ $\vec{B}$
\end{enumerate}

\begin{tcolorbox}[title=นิยามการบวกเวกเตอร์]
ผลบวกของเวกเตอร์คือเวกเตอร์ลัพธ์ที่แทนผลรวมของการเคลื่อนที่หรือผลรวมของปริมาณเวกเตอร์หลายตัว โดยต้องคำนึงถึงทั้งขนาดและทิศทาง
\end{tcolorbox}

\subsection{กฎสลับที่ของการบวกเวกเตอร์}

ถ้าบวก $\vec{A}$ แล้วตามด้วย $\vec{B}$ จะได้เวกเตอร์ลัพธ์เดียวกับบวก $\vec{B}$ แล้วตามด้วย $\vec{A}$

\[
\vec{A}+\vec{B}=\vec{B}+\vec{A}
\]

กฎนี้เรียกว่า กฎสลับที่ของการบวก

\subsection{กฎเปลี่ยนหมู่ของการบวกเวกเตอร์}

ถ้ามีเวกเตอร์สามตัว

\[
\vec{A}+\vec{B}+\vec{C}
\]

จะรวม $\vec{A}$ กับ $\vec{B}$ ก่อน หรือรวม $\vec{B}$ กับ $\vec{C}$ ก่อนก็ได้

\[
\vec{A}+(\vec{B}+\vec{C})=(\vec{A}+\vec{B})+\vec{C}
\]

กฎนี้เรียกว่า กฎเปลี่ยนหมู่ของการบวก

\begin{tcolorbox}[title=ข้อจำกัดสำคัญ]
เวกเตอร์ที่จะนำมาบวกกันต้องเป็นปริมาณชนิดเดียวกัน เช่น แรงบวกกับแรงได้ ความเร็วบวกกับความเร็วได้ แต่แรงบวกกับความเร็วโดยตรงไม่ได้
\end{tcolorbox}

\subsection{การลบเวกเตอร์}

การลบเวกเตอร์ไม่ใช่เรื่องใหม่ เพราะสามารถเปลี่ยนเป็นการบวกเวกเตอร์ตรงข้ามได้

เริ่มจาก

\[
\vec{A}-\vec{B}
\]

เขียนใหม่เป็น

\[
\vec{A}+(-\vec{B})
\]

ดังนั้นการหา $\vec{A}-\vec{B}$ คือการบวก $\vec{A}$ กับเวกเตอร์ที่มีทิศตรงข้ามกับ $\vec{B}$

\begin{tcolorbox}[title=สูตรการลบเวกเตอร์]
\[
\vec{A}-\vec{B}=\vec{A}+(-\vec{B})
\]
\end{tcolorbox}

\subsection{การคูณเวกเตอร์ด้วยจำนวนจริง}

ถ้า $c$ เป็นจำนวนจริง การคูณ

\[
c\vec{A}
\]

มีความหมายดังนี้

\begin{enumerate}
    \item ขนาดเปลี่ยนเป็น $|c|A$
    \item ถ้า $c>0$ ทิศทางเหมือนเดิม
    \item ถ้า $c<0$ ทิศทางกลับตรงข้าม
    \item ถ้า $c=0$ ได้เวกเตอร์ศูนย์
\end{enumerate}

ตัวอย่างเช่น

\[
2\vec{A}
\]

มีขนาดเป็นสองเท่าของ $\vec{A}$ และทิศทางเดิม

\[
-\frac{1}{2}\vec{A}
\]

มีขนาดครึ่งหนึ่งของ $\vec{A}$ แต่ทิศทางตรงข้าม

\begin{examplebox}{ตัวอย่าง: บวกเวกเตอร์ตั้งฉาก}
นักเรียนเดินไปทางตะวันออก $6\,\si{m}$ แล้วเดินไปทางเหนือ $8\,\si{m}$ จงหาขนาดของการกระจัดลัพธ์

เพราะสองทิศตั้งฉากกัน ใช้พีทาโกรัส

\[
R=\sqrt{6^2+8^2}
\]

\[
R=10\,\si{m}
\]

ถ้าต้องการหามุมกับทิศตะวันออก

\[
\tan\theta=\frac{8}{6}
\]

\[
\theta=\tan^{-1}\left(\frac{4}{3}\right)
\]
\end{examplebox}

\begin{exercisebox}{แบบฝึกหัดหน่วยที่ 2}
\begin{enumerate}
    \item นักเรียนเดินไปทางตะวันออก $5\,\si{m}$ แล้วเดินไปทางเหนือ $12\,\si{m}$ จงหาขนาดของการกระจัดลัพธ์
    \item เวกเตอร์ $\vec{A}$ มีขนาด $10\,\si{N}$ ไปทางขวา และ $\vec{B}$ มีขนาด $4\,\si{N}$ ไปทางซ้าย จงหา $\vec{A}+\vec{B}$
    \item ถ้า $\vec{A}$ ไปทางเหนือ และ $\vec{B}$ ไปทางตะวันออก จงอธิบายทิศของ $\vec{A}+\vec{B}$
    \item ถ้า $\vec{A}=3\hat{i}+4\hat{j}$ จงอธิบายความหมายทางเรขาคณิตของ $-\vec{A}$
\end{enumerate}
\end{exercisebox}

% =========================================================
\section{หน่วยที่ 3: องค์ประกอบของเวกเตอร์}

\subsection{ทำไมต้องแตกเวกเตอร์}

การบวกเวกเตอร์ด้วยรูปวาดมีประโยชน์ในการเข้าใจทิศทาง แต่โจทย์สอบต้องการคำตอบที่แม่นยำ วิธีที่แม่นยำกว่าคือการแตกเวกเตอร์เป็นองค์ประกอบตามแกน

แนวคิดคือ แทนเวกเตอร์เอียงหนึ่งตัวด้วยเวกเตอร์สองตัวที่ตั้งฉากกัน

\[
\vec{A}=\vec{A}_x+\vec{A}_y
\]

โดย $\vec{A}_x$ อยู่ตามแกน $x$ และ $\vec{A}_y$ อยู่ตามแกน $y$

\subsection{หาองค์ประกอบจากขนาดและมุม}

ให้เวกเตอร์ $\vec{A}$ มีขนาด $A$ และทำมุม $\theta$ กับแกน $+x$

เริ่มจากสามเหลี่ยมมุมฉาก

ด้านประชิดมุมคือองค์ประกอบตามแกน $x$

\[
\cos\theta=\frac{A_x}{A}
\]

จัดรูปได้

\[
A_x=A\cos\theta
\]

ด้านตรงข้ามมุมคือองค์ประกอบตามแกน $y$

\[
\sin\theta=\frac{A_y}{A}
\]

จัดรูปได้

\[
A_y=A\sin\theta
\]

\begin{tcolorbox}[title=สูตรแตกเวกเตอร์ เมื่อมุมวัดจากแกน $+x$]
\[
A_x=A\cos\theta
\]

\[
A_y=A\sin\theta
\]
\end{tcolorbox}

\begin{tcolorbox}[title=ระวัง]
สูตร $A_x=A\cos\theta$ และ $A_y=A\sin\theta$ ใช้ตรง ๆ เมื่อ $\theta$ วัดจากแกน $+x$ เท่านั้น ถ้ามุมวัดจากแกนอื่น ต้องดูว่าองค์ประกอบใดเป็นด้านประชิดและด้านตรงข้ามมุม
\end{tcolorbox}

\subsection{หาเวกเตอร์จากองค์ประกอบ}

ถ้ารู้ว่าเวกเตอร์มีองค์ประกอบ $A_x$ และ $A_y$ เราหาขนาดของเวกเตอร์จากพีทาโกรัส

\[
A^2=A_x^2+A_y^2
\]

ดังนั้น

\[
A=\sqrt{A_x^2+A_y^2}
\]

และหาทิศทางจาก

\[
\tan\theta=\frac{A_y}{A_x}
\]

ดังนั้น

\[
\theta=\tan^{-1}\left(\frac{A_y}{A_x}\right)
\]

\begin{tcolorbox}[title=สูตรหาขนาดและทิศจากองค์ประกอบ]
\[
A=\sqrt{A_x^2+A_y^2}
\]

\[
\tan\theta=\frac{A_y}{A_x}
\]
\end{tcolorbox}

\subsection{เครื่องหมายขององค์ประกอบ}

เครื่องหมายของ $A_x$ และ $A_y$ บอกทิศทางขององค์ประกอบ

\begin{center}
\begin{tabular}{|c|c|c|}
\hline
\textbf{ควอดแรนต์} & \textbf{เครื่องหมายของ $A_x$} & \textbf{เครื่องหมายของ $A_y$} \\
\hline
I & $+$ & $+$ \\
\hline
II & $-$ & $+$ \\
\hline
III & $-$ & $-$ \\
\hline
IV & $+$ & $-$ \\
\hline
\end{tabular}
\end{center}

\begin{examplebox}{ตัวอย่าง: แตกแรงเป็นองค์ประกอบ}
แรง $\vec{F}$ มีขนาด $50\,\si{N}$ ทำมุม $37^\circ$ กับแกน $+x$ จงหา $F_x$ และ $F_y$

เริ่มจากสูตรเมื่อมุมวัดจากแกน $+x$

\[
F_x=F\cos\theta
\]

\[
F_y=F\sin\theta
\]

แทนค่า

\[
F_x=50\cos37^\circ
\]

\[
F_y=50\sin37^\circ
\]

ถ้าใช้ค่าประมาณ $\cos37^\circ\approx \frac{4}{5}$ และ $\sin37^\circ\approx \frac{3}{5}$

\[
F_x=50\left(\frac{4}{5}\right)=40\,\si{N}
\]

\[
F_y=50\left(\frac{3}{5}\right)=30\,\si{N}
\]
\end{examplebox}

\begin{examplebox}{ตัวอย่าง: มุมวัดจากแกน $y$}
แรง $\vec{F}$ มีขนาด $20\,\si{N}$ ทำมุม $30^\circ$ กับแกน $+y$ ไปทางขวา จงหา $F_x$ และ $F_y$

ครั้งนี้มุมไม่ได้วัดจากแกน $+x$ จึงต้องดูสามเหลี่ยม

องค์ประกอบตามแกน $y$ เป็นด้านประชิดมุม

\[
F_y=F\cos30^\circ
\]

องค์ประกอบตามแกน $x$ เป็นด้านตรงข้ามมุม

\[
F_x=F\sin30^\circ
\]

ดังนั้น

\[
F_x=20\sin30^\circ=10\,\si{N}
\]

\[
F_y=20\cos30^\circ=10\sqrt{3}\,\si{N}
\]
\end{examplebox}

\begin{exercisebox}{แบบฝึกหัดหน่วยที่ 3}
\begin{enumerate}
    \item เวกเตอร์ขนาด $10$ ทำมุม $30^\circ$ กับแกน $+x$ จงหา $A_x$ และ $A_y$
    \item เวกเตอร์หนึ่งมี $A_x=6$ และ $A_y=8$ จงหาขนาดของเวกเตอร์
    \item เวกเตอร์หนึ่งมี $A_x=-3$ และ $A_y=4$ จงหาขนาดและระบุควอดแรนต์
    \item แรง $100\,\si{N}$ ทำมุม $60^\circ$ กับแนวดิ่งไปทางขวา จงหาองค์ประกอบตามแนวนอนและแนวดิ่ง
    \item เวกเตอร์หนึ่งมีขนาด $20$ และชี้ลงทางขวาทำมุม $30^\circ$ ใต้แกน $+x$ จงหาองค์ประกอบ $x$ และ $y$
\end{enumerate}
\end{exercisebox}

% =========================================================
\section{หน่วยที่ 4: เวกเตอร์หน่วย}

\subsection{ทำไมต้องมีเวกเตอร์หน่วย}

เมื่อเราแตกเวกเตอร์เป็นองค์ประกอบแล้ว เราต้องการเขียนให้ชัดว่าแต่ละองค์ประกอบชี้ไปทางแกนใด จึงใช้เวกเตอร์หน่วย

เวกเตอร์หน่วยคือเวกเตอร์ที่มีขนาดเท่ากับ $1$ และใช้บอกทิศทางเท่านั้น

\[
|\hat{i}|=|\hat{j}|=|\hat{k}|=1
\]

โดย

\[
\hat{i} = \text{เวกเตอร์หน่วยตามแกน } +x
\]

\[
\hat{j} = \text{เวกเตอร์หน่วยตามแกน } +y
\]

\[
\hat{k} = \text{เวกเตอร์หน่วยตามแกน } +z
\]

\begin{tcolorbox}[title=ความหมายของเวกเตอร์หน่วย]
เวกเตอร์หน่วยไม่มีหน้าที่บอกขนาดของปริมาณทางฟิสิกส์ แต่ใช้บอกทิศทางขององค์ประกอบ
\end{tcolorbox}

\subsection{เขียนเวกเตอร์ในรูปเวกเตอร์หน่วย}

ถ้าเวกเตอร์ $\vec{A}$ มีองค์ประกอบ

\[
A_x
\]

ตามแกน $x$ และ

\[
A_y
\]

ตามแกน $y$ เราเขียนได้ว่า

\[
\vec{A}=A_x\hat{i}+A_y\hat{j}
\]

ในสามมิติ เขียนเป็น

\[
\vec{A}=A_x\hat{i}+A_y\hat{j}+A_z\hat{k}
\]

\begin{examplebox}{ตัวอย่าง}
เวกเตอร์หนึ่งมีองค์ประกอบ $x$ เท่ากับ $3$ และองค์ประกอบ $y$ เท่ากับ $-4$

จึงเขียนได้ว่า

\[
\vec{A}=3\hat{i}-4\hat{j}
\]

ขนาดของเวกเตอร์คือ

\[
A=\sqrt{3^2+(-4)^2}=5
\]

เวกเตอร์นี้ชี้ไปควอดแรนต์ที่ IV เพราะ $A_x>0$ และ $A_y<0$
\end{examplebox}

\subsection{เวกเตอร์ตำแหน่ง}

ถ้าจุด $P$ อยู่ที่ตำแหน่ง

\[
(x,y)
\]

เวกเตอร์จากจุดกำเนิดไปยังจุด $P$ เรียกว่า เวกเตอร์ตำแหน่ง

\[
\vec{r}=x\hat{i}+y\hat{j}
\]

ในสามมิติ

\[
\vec{r}=x\hat{i}+y\hat{j}+z\hat{k}
\]

\begin{tcolorbox}[title=เชื่อมกับบทการเคลื่อนที่]
ในบทต่อไป ตำแหน่งของอนุภาคในสองมิติจะเขียนเป็นเวกเตอร์ตำแหน่ง $\vec{r}(t)$ และความเร็วจะได้จากอัตราการเปลี่ยนของเวกเตอร์ตำแหน่งตามเวลา
\end{tcolorbox}

\begin{exercisebox}{แบบฝึกหัดหน่วยที่ 4}
\begin{enumerate}
    \item จงเขียนเวกเตอร์ที่มี $A_x=5$ และ $A_y=12$ ในรูปเวกเตอร์หน่วย
    \item จงหาขนาดของ $\vec{A}=8\hat{i}-6\hat{j}$
    \item จุด $P$ อยู่ที่ $(4,-7)$ จงเขียนเวกเตอร์ตำแหน่ง $\vec{r}$
    \item เวกเตอร์ $\vec{B}=-3\hat{i}-4\hat{j}$ อยู่ควอดแรนต์ใด และมีขนาดเท่าไร
    \item เวกเตอร์หนึ่งในสามมิติคือ $\vec{C}=2\hat{i}-3\hat{j}+6\hat{k}$ จงหาขนาดของเวกเตอร์นี้
\end{enumerate}
\end{exercisebox}

% =========================================================
\section{หน่วยที่ 5: การบวกเวกเตอร์ด้วยองค์ประกอบ}

\subsection{แนวคิดหลัก}

การบวกเวกเตอร์ด้วยรูปวาดช่วยให้เห็นภาพ แต่ถ้าต้องการคำตอบแม่นยำ ให้แตกทุกเวกเตอร์เป็นองค์ประกอบก่อน แล้วบวกทีละแกน

ให้

\[
\vec{A}=A_x\hat{i}+A_y\hat{j}
\]

และ

\[
\vec{B}=B_x\hat{i}+B_y\hat{j}
\]

ต้องการหา

\[
\vec{R}=\vec{A}+\vec{B}
\]

แทนรูปองค์ประกอบลงไป

\[
\vec{R}=(A_x\hat{i}+A_y\hat{j})+(B_x\hat{i}+B_y\hat{j})
\]

จัดกลุ่มแกนเดียวกัน

\[
\vec{R}=(A_x+B_x)\hat{i}+(A_y+B_y)\hat{j}
\]

ดังนั้นองค์ประกอบของเวกเตอร์ลัพธ์คือ

\[
R_x=A_x+B_x
\]

\[
R_y=A_y+B_y
\]

\begin{tcolorbox}[title=สูตรบวกเวกเตอร์ด้วยองค์ประกอบ]
\[
R_x=A_x+B_x
\]

\[
R_y=A_y+B_y
\]

\[
\vec{R}=R_x\hat{i}+R_y\hat{j}
\]
\end{tcolorbox}

\subsection{หาขนาดและทิศของเวกเตอร์ลัพธ์}

หลังจากได้ $R_x$ และ $R_y$ แล้ว ใช้พีทาโกรัส

\[
R=\sqrt{R_x^2+R_y^2}
\]

และใช้แทนเจนต์หามุม

\[
\tan\theta=\frac{R_y}{R_x}
\]

ดังนั้น

\[
\theta=\tan^{-1}\left(\frac{R_y}{R_x}\right)
\]

แต่ต้องตรวจควอดแรนต์จากเครื่องหมายของ $R_x$ และ $R_y$

\begin{examplebox}{ตัวอย่าง: บวกเวกเตอร์ด้วยองค์ประกอบ}
กำหนดให้

\[
\vec{A}=3\hat{i}+4\hat{j}
\]

และ

\[
\vec{B}=5\hat{i}-2\hat{j}
\]

ต้องการหา

\[
\vec{R}=\vec{A}+\vec{B}
\]

บวกองค์ประกอบตามแกน $x$

\[
R_x=3+5=8
\]

บวกองค์ประกอบตามแกน $y$

\[
R_y=4+(-2)=2
\]

ดังนั้น

\[
\vec{R}=8\hat{i}+2\hat{j}
\]

ขนาดคือ

\[
R=\sqrt{8^2+2^2}
\]

\[
R=\sqrt{68}=2\sqrt{17}
\]

มุมกับแกน $+x$ หาได้จาก

\[
\tan\theta=\frac{2}{8}=\frac{1}{4}
\]

\[
\theta=\tan^{-1}\left(\frac{1}{4}\right)
\]
\end{examplebox}

\subsection{การลบเวกเตอร์ด้วยองค์ประกอบ}

ให้

\[
\vec{D}=\vec{A}-\vec{B}
\]

จากนิยามการลบเวกเตอร์

\[
\vec{A}-\vec{B}=\vec{A}+(-\vec{B})
\]

ในรูปองค์ประกอบจึงได้

\[
D_x=A_x-B_x
\]

\[
D_y=A_y-B_y
\]

ดังนั้น

\[
\vec{D}=(A_x-B_x)\hat{i}+(A_y-B_y)\hat{j}
\]

\begin{examplebox}{ตัวอย่าง: ลบเวกเตอร์}
กำหนดให้

\[
\vec{A}=7\hat{i}+3\hat{j}
\]

และ

\[
\vec{B}=2\hat{i}+9\hat{j}
\]

หา

\[
\vec{D}=\vec{A}-\vec{B}
\]

องค์ประกอบตามแกน $x$

\[
D_x=7-2=5
\]

องค์ประกอบตามแกน $y$

\[
D_y=3-9=-6
\]

ดังนั้น

\[
\vec{D}=5\hat{i}-6\hat{j}
\]
\end{examplebox}

\begin{exercisebox}{แบบฝึกหัดหน่วยที่ 5}
\begin{enumerate}
    \item ให้ $\vec{A}=2\hat{i}+5\hat{j}$ และ $\vec{B}=4\hat{i}-3\hat{j}$ จงหา $\vec{A}+\vec{B}$
    \item ให้ $\vec{A}=6\hat{i}-2\hat{j}$ และ $\vec{B}=-3\hat{i}+7\hat{j}$ จงหา $\vec{A}-\vec{B}$
    \item เวกเตอร์สองตัวคือ $\vec{A}=10\hat{i}$ และ $\vec{B}=10\hat{j}$ จงหาขนาดของ $\vec{A}+\vec{B}$
    \item เวกเตอร์ลัพธ์มี $R_x=-8$ และ $R_y=6$ จงหาขนาดและควอดแรนต์ของเวกเตอร์ลัพธ์
    \item เวกเตอร์ $\vec{A}$ มีขนาด $20$ ทำมุม $30^\circ$ กับแกน $+x$ และ $\vec{B}$ มีขนาด $10$ ชี้ไปทางแกน $-y$ จงหาองค์ประกอบของ $\vec{A}+\vec{B}$
\end{enumerate}
\end{exercisebox}

% =========================================================
\section{การประยุกต์กับโจทย์ สอวน.}

\subsection{แนวคิดการแก้โจทย์เวกเตอร์แบบ สอวน.}

โจทย์ สอวน. มักไม่ได้ถามเวกเตอร์ตรง ๆ ว่า ``จงบวกเวกเตอร์'' แต่ซ่อนเวกเตอร์ไว้ในบริบทอื่น เช่น

\[
\text{แรง}
\]

\[
\text{โมเมนตัม}
\]

\[
\text{สนามไฟฟ้า}
\]

\[
\text{ความเร็วโพรเจคไทล์}
\]

\[
\text{การชนหรือการระเบิด}
\]

หลักคิดคือ ให้แยกปัญหาออกเป็นแกน $x$ และ $y$ เสมอ

\begin{tcolorbox}[title=ขั้นตอนแก้โจทย์เวกเตอร์]
\begin{enumerate}
    \item วาดแกนและเลือกทิศบวก
    \item วาดเวกเตอร์ทุกตัว
    \item แตกเวกเตอร์เป็นองค์ประกอบตามแกน
    \item ตั้งสมการแยกแกน $x$ และแกน $y$
    \item รวมคำตอบกลับเป็นขนาดและทิศทางเมื่อต้องการ
\end{enumerate}
\end{tcolorbox}

\subsection{ตัวอย่างแนว สอวน.: เวกเตอร์ความเร็วหลังการระเบิด}
\vspace{1.5em}
\begin{center}
    \includegraphics[width=1\linewidth]{img/posn66q3.png}
\end{center}
\begin{examplebox}{เฉลยตัวอย่างโจทย์ สอวน.}

มวล $m$ เคลื่อนที่ด้วยความเร็ว $u$ ตามแกน $+x$ ต่อมาระเบิดออกเป็นสองส่วน ส่วนแรกมีมวล $\alpha m$ และเคลื่อนที่ด้วยความเร็ว $v$ ตามแกน $+y$ จงพิจารณาทิศทางของส่วนที่สอง

ก่อนระเบิด โมเมนตัมรวมอยู่ตามแกน $x$ เท่านั้น

\[
\vec{p}_{\text{before}}=mu\hat{i}
\]

หลังระเบิด ส่วนแรกมีโมเมนตัม

\[
\vec{p}_1=\alpha mv\hat{j}
\]

ให้ส่วนที่สองมีมวล

\[
(1-\alpha)m
\]

และมีความเร็ว

\[
\vec{v}_2=v_{2x}\hat{i}+v_{2y}\hat{j}
\]

โมเมนตัมของส่วนที่สองคือ

\[
\vec{p}_2=(1-\alpha)m(v_{2x}\hat{i}+v_{2y}\hat{j})
\]

ใช้การอนุรักษ์โมเมนตัม

\[
\vec{p}_{\text{before}}=\vec{p}_1+\vec{p}_2
\]

แยกแกน $x$

\[
mu=(1-\alpha)mv_{2x}
\]

ดังนั้น

\[
v_{2x}=\frac{u}{1-\alpha}
\]

แยกแกน $y$

\[
0=\alpha mv+(1-\alpha)mv_{2y}
\]

ดังนั้น

\[
v_{2y}=-\frac{\alpha v}{1-\alpha}
\]

ถ้าให้ $\theta$ เป็นมุมที่ส่วนที่สองทำกับแกน $+x$ ลงด้านล่าง จะได้

\[
\tan\theta=\frac{|v_{2y}|}{v_{2x}}
\]

\[
\tan\theta=\frac{\alpha v}{u}
\]

\[
\theta=\tan^{-1}\left(\frac{\alpha v}{u}\right)
\]
\end{examplebox}



\subsection{ตัวอย่างแนว สอวน.: โพรเจคไทล์กับพื้นเอียง}

\vspace{1.5em}
\begin{center}
    \includegraphics[width=1\linewidth]{img/posn67q2.png}
\end{center}

\begin{examplebox}{แนวคิดการแก้โจทย์}

ในโจทย์โพรเจคไทล์ ความเร็วต้นมักเป็นเวกเตอร์ที่ต้องแตกเป็นสองแกน

ถ้าวัตถุถูกยิงด้วยความเร็วต้น $u$ ทำมุม $\theta$ กับแกน $+x$

เริ่มจากการแตกเวกเตอร์ความเร็วต้น

\[
u_x=u\cos\theta
\]

\[
u_y=u\sin\theta
\]

หลังจากนั้น แกน $x$ และแกน $y$ แยกคิดได้

แกน $x$ ไม่มีความเร่งถ้าไม่คิดแรงต้านอากาศ

\[
x=u_xt
\]

แกน $y$ มีความเร่งโน้มถ่วง

\[
y=u_yt-\frac{1}{2}gt^2
\]

\end{examplebox}

\begin{exercisebox}{แบบฝึกหัดแนว สอวน.}
\begin{enumerate}
    \item มวลหนึ่งเคลื่อนที่ด้วยความเร็ว $u$ ตามแกน $+x$ แล้วแตกออกเป็นสองชิ้น ชิ้นแรกเคลื่อนที่ตามแกน $+y$ จงอธิบายว่าทำไมชิ้นที่สองต้องมีองค์ประกอบความเร็วตามแกน $-y$
    \item จุดประจุ $+q$ สองตัวอยู่ที่ $(0,a)$ และ $(0,-a)$ จงพิจารณาทิศของสนามไฟฟ้าลัพธ์ที่จุด $(x,0)$ โดยไม่ต้องคำนวณขนาดเต็ม
    \item วัตถุถูกยิงด้วยความเร็วต้น $u$ ทำมุม $30^\circ$ กับแนวระดับ จงเขียน $u_x$ และ $u_y$
    \item แรงสองแรงขนาด $F$ เท่ากันทำมุม $\theta$ กับแนวดิ่งคนละด้าน จงหาองค์ประกอบแนวดิ่งของแรงลัพธ์
    \item เวกเตอร์สนามไฟฟ้าสองตัวมีขนาดเท่ากันและสมมาตรกับแกน $x$ จงอธิบายว่าองค์ประกอบตามแกนใดหักล้างกัน และองค์ประกอบตามแกนใดรวมกัน
\end{enumerate}
\end{exercisebox}

% =========================================================
\section{คณิตศาสตร์ท้ายบท: เวกเตอร์ที่เป็นฟังก์ชันของเวลา}

แม้บทนี้เน้นพีชคณิตและตรีโกณมิติ แต่ในบทถัดไปเราจะใช้เวกเตอร์กับการเคลื่อนที่สองมิติ ดังนั้นควรเห็นภาพล่วงหน้าว่าเวกเตอร์สามารถเปลี่ยนตามเวลาได้

ถ้าตำแหน่งของอนุภาคในสองมิติเขียนเป็น

\[
\vec{r}(t)=x(t)\hat{i}+y(t)\hat{j}
\]

ความเร็วคืออัตราการเปลี่ยนของตำแหน่งตามเวลา

\[
\vec{v}(t)=\frac{d\vec{r}}{dt}
\]

เพราะ $\hat{i}$ และ $\hat{j}$ เป็นทิศคงที่ในระบบพิกัดฉาก เราหาอนุพันธ์ทีละองค์ประกอบได้

\[
\vec{v}(t)=\frac{dx}{dt}\hat{i}+\frac{dy}{dt}\hat{j}
\]

ดังนั้น

\[
\vec{v}(t)=v_x(t)\hat{i}+v_y(t)\hat{j}
\]

โดย

\[
v_x(t)=\frac{dx}{dt}
\]

\[
v_y(t)=\frac{dy}{dt}
\]

ในทำนองเดียวกัน ความเร่งคือ

\[
\vec{a}(t)=\frac{d\vec{v}}{dt}
\]

\[
\vec{a}(t)=\frac{dv_x}{dt}\hat{i}+\frac{dv_y}{dt}\hat{j}
\]

ดังนั้น

\[
\vec{a}(t)=a_x(t)\hat{i}+a_y(t)\hat{j}
\]

\begin{examplebox}{ตัวอย่างเตรียมบทถัดไป}
กำหนดให้

\[
\vec{r}(t)=(3t^2)\hat{i}+(4t)\hat{j}
\]

ต้องการหาความเร็ว

\[
\vec{v}(t)=\frac{d\vec{r}}{dt}
\]

หาอนุพันธ์ทีละแกน

\[
\vec{v}(t)=6t\hat{i}+4\hat{j}
\]

ต้องการหาความเร่ง

\[
\vec{a}(t)=\frac{d\vec{v}}{dt}
\]

ดังนั้น

\[
\vec{a}(t)=6\hat{i}+0\hat{j}
\]

หรือ

\[
\vec{a}(t)=6\hat{i}
\]
\end{examplebox}

% =========================================================
\section{ชุดโจทย์รวมท้ายบท: ระดับผสมสำหรับ สอวน. ค่าย 1}

\subsection*{ระดับพื้นฐาน}

\begin{enumerate}
    \item จงแปลงเวกเตอร์ขนาด $10$ ทำมุม $60^\circ$ กับแกน $+x$ ให้อยู่ในรูปองค์ประกอบ
    \item จงหาขนาดของเวกเตอร์ $\vec{A}=5\hat{i}+12\hat{j}$
    \item ให้ $\vec{A}=3\hat{i}+4\hat{j}$ และ $\vec{B}=2\hat{i}-7\hat{j}$ จงหา $\vec{A}+\vec{B}$
    \item จุด $P$ มีพิกัด $(-6,8)$ จงหา $r$ และบอกควอดแรนต์
\end{enumerate}

\subsection*{ระดับกลาง}

\begin{enumerate}
    \setcounter{enumi}{4}
    \item เวกเตอร์ $\vec{A}$ มีขนาด $20$ ทำมุม $30^\circ$ กับแกน $+x$ และเวกเตอร์ $\vec{B}$ มีขนาด $10$ ทำมุม $120^\circ$ กับแกน $+x$ จงหาองค์ประกอบของ $\vec{A}+\vec{B}$
    \item แรงสองแรงขนาด $F$ เท่ากัน ทำมุม $\theta$ กับแนวดิ่งคนละด้าน จงหาแรงลัพธ์
    \item เวกเตอร์หนึ่งมีขนาด $A$ และมี $A_x=-\frac{A}{2}$ จงหามุมที่เวกเตอร์ทำกับแกน $+x$
    \item ให้ $\vec{A}=a\hat{i}+2a\hat{j}$ และ $\vec{B}=3a\hat{i}-a\hat{j}$ จงหาขนาดของ $\vec{A}-\vec{B}$
\end{enumerate}

\subsection*{ระดับยาก}

\begin{enumerate}
    \setcounter{enumi}{8}
    \item เวกเตอร์สามตัวมีผลรวมเป็นศูนย์ ถ้า
    \[
    \vec{A}=4\hat{i}+3\hat{j}
    \]
    และ
    \[
    \vec{B}=-2\hat{i}+5\hat{j}
    \]
    จงหา $\vec{C}$
    
    \item เวกเตอร์ $\vec{A}$ มีขนาด $A$ ทำมุม $\theta$ กับแกน $+x$ เวกเตอร์ $\vec{B}$ มีขนาด $B$ ทำมุม $-\theta$ กับแกน $+x$ จงหาองค์ประกอบของ $\vec{A}+\vec{B}$
    
    \item มวลหนึ่งมีโมเมนตัมเริ่มต้น $mu\hat{i}$ แล้วแตกออกเป็นสองชิ้น ชิ้นแรกมีมวล $\alpha m$ และมีความเร็ว $v\hat{j}$ จงหาทิศของความเร็วชิ้นที่สองในรูปของ $\alpha,u,v$
    
    \item วัตถุถูกยิงด้วยความเร็วต้น $u$ ทำมุม $\theta$ กับแนวระดับ จงเขียนเวกเตอร์ความเร็วต้นในรูปเวกเตอร์หน่วย
    
    \item เวกเตอร์ตำแหน่งของอนุภาคคือ
    \[
    \vec{r}(t)=(2t^2-1)\hat{i}+(3t+4)\hat{j}
    \]
    จงหา $\vec{v}(t)$ และ $\vec{a}(t)$
    
    \item เวกเตอร์สองตัวมีขนาดเท่ากันคือ $A$ และทำมุมระหว่างกัน $60^\circ$ จงหาขนาดของผลบวก โดยใช้การแตกองค์ประกอบ
\end{enumerate}

% =========================================================
\section{เฉลยย่อท้ายบท}

\begin{enumerate}
    \item
    \[
    A_x=10\cos60^\circ=5
    \]
    \[
    A_y=10\sin60^\circ=5\sqrt{3}
    \]
    
    \item
    \[
    A=\sqrt{5^2+12^2}=13
    \]
    
    \item
    \[
    \vec{A}+\vec{B}=(3+2)\hat{i}+(4-7)\hat{j}
    \]
    \[
    \vec{A}+\vec{B}=5\hat{i}-3\hat{j}
    \]
    
    \item
    \[
    r=\sqrt{(-6)^2+8^2}=10
    \]
    จุดอยู่ควอดแรนต์ที่ II
    
    \item
    \[
    A_x=20\cos30^\circ=10\sqrt{3}
    \]
    \[
    A_y=20\sin30^\circ=10
    \]
    \[
    B_x=10\cos120^\circ=-5
    \]
    \[
    B_y=10\sin120^\circ=5\sqrt{3}
    \]
    ดังนั้น
    \[
    R_x=10\sqrt{3}-5
    \]
    \[
    R_y=10+5\sqrt{3}
    \]
    
    \item
    องค์ประกอบแนวนอนหักล้างกัน และองค์ประกอบแนวดิ่งรวมกัน
    \[
    R=2F\cos\theta
    \]
    ทิศตามแนวดิ่ง
    
    \item
    \[
    \cos\theta=\frac{A_x}{A}=-\frac{1}{2}
    \]
    ดังนั้น
    \[
    \theta=120^\circ
    \]
    หรือถ้าอยู่ควอดแรนต์ III อาจเป็น $240^\circ$ ต้องใช้ข้อมูลของ $A_y$ เพิ่ม
    
    \item
    \[
    \vec{A}-\vec{B}=(a-3a)\hat{i}+(2a-(-a))\hat{j}
    \]
    \[
    \vec{A}-\vec{B}=-2a\hat{i}+3a\hat{j}
    \]
    \[
    |\vec{A}-\vec{B}|=\sqrt{(-2a)^2+(3a)^2}=\sqrt{13}a
    \]
    
    \item
    \[
    \vec{A}+\vec{B}+\vec{C}=0
    \]
    \[
    \vec{C}=-(\vec{A}+\vec{B})
    \]
    \[
    \vec{A}+\vec{B}=2\hat{i}+8\hat{j}
    \]
    \[
    \vec{C}=-2\hat{i}-8\hat{j}
    \]
    
    \item
    \[
    A_x=A\cos\theta,\quad A_y=A\sin\theta
    \]
    \[
    B_x=B\cos\theta,\quad B_y=-B\sin\theta
    \]
    ดังนั้น
    \[
    R_x=(A+B)\cos\theta
    \]
    \[
    R_y=(A-B)\sin\theta
    \]
    
    \item
    ใช้การอนุรักษ์โมเมนตัม
    \[
    \tan\theta=\frac{\alpha v}{u}
    \]
    โดย $\theta$ เป็นมุมที่ชิ้นที่สองทำกับแกน $+x$ ลงด้านล่าง
    
    \item
    \[
    \vec{u}=u\cos\theta\hat{i}+u\sin\theta\hat{j}
    \]
    
    \item
    \[
    \vec{v}(t)=\frac{d\vec{r}}{dt}=4t\hat{i}+3\hat{j}
    \]
    \[
    \vec{a}(t)=\frac{d\vec{v}}{dt}=4\hat{i}
    \]
    
    \item
    เลือกให้เวกเตอร์แรกอยู่ตามแกน $x$
    \[
    \vec{A}=A\hat{i}
    \]
    เวกเตอร์ที่สองทำมุม $60^\circ$
    \[
    \vec{B}=A\cos60^\circ\hat{i}+A\sin60^\circ\hat{j}
    \]
    ดังนั้น
    \[
    R_x=A+\frac{A}{2}=\frac{3A}{2}
    \]
    \[
    R_y=\frac{\sqrt{3}A}{2}
    \]
    \[
    R=\sqrt{\left(\frac{3A}{2}\right)^2+\left(\frac{\sqrt{3}A}{2}\right)^2}
    \]
    \[
    R=\sqrt{\frac{9A^2}{4}+\frac{3A^2}{4}}
    \]
    \[
    R=\sqrt{3A^2}=A\sqrt{3}
    \]
\end{enumerate}

% =========================================================
\section{สรุปสูตรสำคัญ}

\begin{tcolorbox}[title=สรุปบทที่ 3: เวกเตอร์]
พิกัดฉากและพิกัดเชิงขั้ว

\[
x=r\cos\theta
\]

\[
y=r\sin\theta
\]

\[
r=\sqrt{x^2+y^2}
\]

\[
\tan\theta=\frac{y}{x}
\]

องค์ประกอบของเวกเตอร์

\[
A_x=A\cos\theta
\]

\[
A_y=A\sin\theta
\]

\[
A=\sqrt{A_x^2+A_y^2}
\]

\[
\tan\theta=\frac{A_y}{A_x}
\]

เวกเตอร์หน่วย

\[
\vec{A}=A_x\hat{i}+A_y\hat{j}
\]

\[
\vec{A}=A_x\hat{i}+A_y\hat{j}+A_z\hat{k}
\]

การบวกเวกเตอร์

\[
\vec{R}=\vec{A}+\vec{B}
\]

\[
R_x=A_x+B_x
\]

\[
R_y=A_y+B_y
\]

\[
R=\sqrt{R_x^2+R_y^2}
\]

\[
\tan\theta=\frac{R_y}{R_x}
\]
\end{tcolorbox}

\begin{tcolorbox}[title=สรุปบทที่ 3: เวกเตอร์ (ต่อ)]
\
การลบเวกเตอร์

\[
\vec{A}-\vec{B}=\vec{A}+(-\vec{B})
\]

\[
D_x=A_x-B_x
\]

\[
D_y=A_y-B_y
\]
\end{tcolorbox}
%----------------------------------------------------------------------------------------
%	INTRODUCTION
%----------------------------------------------------------------------------------------

% \section*{Introduction} % Unnumbered section

% Lorem ipsum dolor sit amet, consectetur adipiscing elit. Praesent porttitor arcu luctus, imperdiet urna iaculis, mattis eros. Pellentesque iaculis odio vel nisl ullamcorper, nec faucibus ipsum molestie. Sed dictum nisl non aliquet porttitor. Etiam vulputate arcu dignissim, finibus sem et, viverra nisl. Aenean luctus congue massa, ut laoreet metus ornare in. Nunc fermentum nisi imperdiet lectus tincidunt vestibulum at ac elit. Nulla mattis nisl eu malesuada suscipit.

% % Math equation/formula
% \begin{equation}
% 	I = \int_{a}^{b} f(x) \; \text{d}x.
% \end{equation}

% Aliquam arcu turpis, ultrices sed luctus ac, vehicula id metus. Morbi eu feugiat velit, et tempus augue. Proin ac mattis tortor. Donec tincidunt, ante rhoncus luctus semper, arcu lorem lobortis justo, nec convallis ante quam quis lectus. Aenean tincidunt sodales massa, et hendrerit tellus mattis ac. Sed non pretium nibh. Donec cursus maximus luctus. Vivamus lobortis eros et massa porta porttitor.

% \begin{info} % Information block
% 	This is an interesting piece of information, to which the reader should pay special attention. Fusce varius orci ac magna dapibus porttitor. In tempor leo a neque bibendum sollicitudin. Nulla pretium fermentum nisi, eget sodales magna facilisis eu. Praesent aliquet nulla ut bibendum lacinia. Donec vel mauris vulputate, commodo ligula ut, egestas orci. Suspendisse commodo odio sed hendrerit lobortis. Donec finibus eros erat, vel ornare enim mattis et.
% \end{info}

% %----------------------------------------------------------------------------------------
% %	PROBLEM 1
% %----------------------------------------------------------------------------------------

% \section{Problem title} % Numbered section

% In hac habitasse platea dictumst. Curabitur mattis elit sit amet justo luctus vestibulum. In hac habitasse platea dictumst. Pellentesque lobortis justo enim, a condimentum massa tempor eu. Ut quis nulla a quam pretium eleifend nec eu nisl. Nam cursus porttitor eros, sed luctus ligula convallis quis. Nam convallis, ligula in auctor euismod, ligula mauris fringilla tellus, et egestas mauris odio eget diam. Praesent sodales in ipsum eu dictum.

% %------------------------------------------------

% \subsection{Theoretical viewpoint}

% Maecenas consectetur metus at tellus finibus condimentum. Proin arcu lectus, ultrices non tincidunt et, tincidunt ut quam. Integer luctus posuere est, non maximus ante dignissim quis. Nunc a cursus erat. Curabitur suscipit nibh in tincidunt sagittis. Nam malesuada vestibulum quam id gravida. Proin ut dapibus velit. Vestibulum eget quam quis ipsum semper convallis. Duis consectetur nibh ac diam dignissim, id condimentum enim dictum. Nam aliquet ligula eu magna pellentesque, nec sagittis leo lobortis. Aenean tincidunt dignissim egestas. Morbi efficitur risus ante, id tincidunt odio pulvinar vitae.

% Curabitur tempus hendrerit nulla. Donec faucibus lobortis nibh pharetra sagittis. Sed magna sem, posuere eget sem vitae, finibus consequat libero. Cras aliquet sagittis erat ut semper. Aenean vel enim ipsum. Fusce ut felis at eros sagittis bibendum mollis lobortis libero. Donec laoreet nisl vel risus lacinia elementum non nec lacus. Nullam luctus, nulla volutpat ultricies ultrices, quam massa placerat augue, ut fringilla urna lectus nec nibh. Vestibulum efficitur condimentum orci a semper. Pellentesque ut metus pretium lacus maximus semper. Sed tellus augue, consectetur rhoncus eleifend vel, imperdiet nec turpis. Nulla ligula ante, malesuada quis orci a, ultricies blandit elit.

% % Numbered question, with subquestions in an enumerate environment
% \begin{question}
% 	Quisque ullamcorper placerat ipsum. Cras nibh. Morbi vel justo vitae lacus tincidunt ultrices. Lorem ipsum dolor sit amet, consectetuer adipiscing elit.

% 	% Subquestions numbered with letters
% 	\begin{enumerate}[(a)]
% 		\item Do this.
% 		\item Do that.
% 		\item Do something else.
% 	\end{enumerate}
% \end{question}
	
% %------------------------------------------------

% \subsection{Algorithmic issues}

% In malesuada ullamcorper urna, sed dapibus diam sollicitudin non. Donec elit odio, accumsan ac nisl a, tempor imperdiet eros. Donec porta tortor eu risus consequat, a pharetra tortor tristique. Morbi sit amet laoreet erat. Morbi et luctus diam, quis porta ipsum. Quisque libero dolor, suscipit id facilisis eget, sodales volutpat dolor. Nullam vulputate interdum aliquam. Mauris id convallis erat, ut vehicula neque. Sed auctor nibh et elit fringilla, nec ultricies dui sollicitudin. Vestibulum vestibulum luctus metus venenatis facilisis. Suspendisse iaculis augue at vehicula ornare. Sed vel eros ut velit fermentum porttitor sed sed massa. Fusce venenatis, metus a rutrum sagittis, enim ex maximus velit, id semper nisi velit eu purus.

% \begin{center}
% 	\begin{minipage}{0.5\linewidth} % Adjust the minipage width to accomodate for the length of algorithm lines
% 		\begin{algorithm}[H]
% 			\KwIn{$(a, b)$, two floating-point numbers}  % Algorithm inputs
% 			\KwResult{$(c, d)$, such that $a+b = c + d$} % Algorithm outputs/results
% 			\medskip
% 			\If{$\vert b\vert > \vert a\vert$}{
% 				exchange $a$ and $b$ \;
% 			}
% 			$c \leftarrow a + b$ \;
% 			$z \leftarrow c - a$ \;
% 			$d \leftarrow b - z$ \;
% 			{\bf return} $(c,d)$ \;
% 			\caption{\texttt{FastTwoSum}} % Algorithm name
% 			\label{alg:fastTwoSum}   % optional label to refer to
% 		\end{algorithm}
% 	\end{minipage}
% \end{center}

% Fusce varius orci ac magna dapibus porttitor. In tempor leo a neque bibendum sollicitudin. Nulla pretium fermentum nisi, eget sodales magna facilisis eu. Praesent aliquet nulla ut bibendum lacinia. Donec vel mauris vulputate, commodo ligula ut, egestas orci. Suspendisse commodo odio sed hendrerit lobortis. Donec finibus eros erat, vel ornare enim mattis et.

% % Numbered question, with an optional title
% \begin{question}[\itshape (with optional title)]
% 	In congue risus leo, in gravida enim viverra id. Donec eros mauris, bibendum vel dui at, tempor commodo augue. In vel lobortis lacus. Nam ornare ullamcorper mauris vel molestie. Maecenas vehicula ornare turpis, vitae fringilla orci consectetur vel. Nam pulvinar justo nec neque egestas tristique. Donec ac dolor at libero congue varius sed vitae lectus. Donec et tristique nulla, sit amet scelerisque orci. Maecenas a vestibulum lectus, vitae gravida nulla. Proin eget volutpat orci. Morbi eu aliquet turpis. Vivamus molestie urna quis tempor tristique. Proin hendrerit sem nec tempor sollicitudin.
% \end{question}

% Mauris interdum porttitor fringilla. Proin tincidunt sodales leo at ornare. Donec tempus magna non mauris gravida luctus. Cras vitae arcu vitae mauris eleifend scelerisque. Nam sem sapien, vulputate nec felis eu, blandit convallis risus. Pellentesque sollicitudin venenatis tincidunt. In et ipsum libero. Nullam tempor ligula a massa convallis pellentesque.

% %----------------------------------------------------------------------------------------
% %	PROBLEM 2
% %----------------------------------------------------------------------------------------

% \section{Implementation}

% Proin lobortis efficitur dictum. Pellentesque vitae pharetra eros, quis dignissim magna. Sed tellus leo, semper non vestibulum vel, tincidunt eu mi. Aenean pretium ut velit sed facilisis. Ut placerat urna facilisis dolor suscipit vehicula. Ut ut auctor nunc. Nulla non massa eros. Proin rhoncus arcu odio, eu lobortis metus sollicitudin eu. Duis maximus ex dui, id bibendum diam dignissim id. Aliquam quis lorem lorem. Phasellus sagittis aliquet dolor, vulputate cursus dolor convallis vel. Suspendisse eu tellus feugiat, bibendum lectus quis, fermentum nunc. Nunc euismod condimentum magna nec bibendum. Curabitur elementum nibh eu sem cursus, eu aliquam leo rutrum. Sed bibendum augue sit amet pharetra ullamcorper. Aenean congue sit amet tortor vitae feugiat.

% In congue risus leo, in gravida enim viverra id. Donec eros mauris, bibendum vel dui at, tempor commodo augue. In vel lobortis lacus. Nam ornare ullamcorper mauris vel molestie. Maecenas vehicula ornare turpis, vitae fringilla orci consectetur vel. Nam pulvinar justo nec neque egestas tristique. Donec ac dolor at libero congue varius sed vitae lectus. Donec et tristique nulla, sit amet scelerisque orci. Maecenas a vestibulum lectus, vitae gravida nulla. Proin eget volutpat orci. Morbi eu aliquet turpis. Vivamus molestie urna quis tempor tristique. Proin hendrerit sem nec tempor sollicitudin.

% % File contents
% \begin{file}[hello.py]
% \begin{lstlisting}[language=Python]
% #! /usr/bin/python

% import sys
% sys.stdout.write("Hello World!\n")
% \end{lstlisting}
% \end{file}

% Fusce eleifend porttitor arcu, id accumsan elit pharetra eget. Mauris luctus velit sit amet est sodales rhoncus. Donec cursus suscipit justo, sed tristique ipsum fermentum nec. Ut tortor ex, ullamcorper varius congue in, efficitur a tellus. Vivamus ut rutrum nisi. Phasellus sit amet enim efficitur, aliquam nulla id, lacinia mauris. Quisque viverra libero ac magna maximus efficitur. Interdum et malesuada fames ac ante ipsum primis in faucibus. Vestibulum mollis eros in tellus fermentum, vitae tristique justo finibus. Sed quis vehicula nibh. Etiam nulla justo, pellentesque id sapien at, semper aliquam arcu. Integer at commodo arcu. Quisque dapibus ut lacus eget vulputate.

% % Command-line "screenshot"
% \begin{commandline}
% 	\begin{verbatim}
% 		$ chmod +x hello.py
% 		$ ./hello.py

% 		Hello World!
% 	\end{verbatim}
% \end{commandline}

% Vestibulum sodales orci a nisi interdum tristique. In dictum vehicula dui, eget bibendum purus elementum eu. Pellentesque lobortis mattis mauris, non feugiat dolor vulputate a. Cras porttitor dapibus lacus at pulvinar. Praesent eu nunc et libero porttitor malesuada tempus quis massa. Aenean cursus ipsum a velit ultricies sagittis. Sed non leo ullamcorper, suscipit massa ut, pulvinar erat. Aliquam erat volutpat. Nulla non lacus vitae mi placerat tincidunt et ac diam. Aliquam tincidunt augue sem, ut vestibulum est volutpat eget. Suspendisse potenti. Integer condimentum, risus nec maximus elementum, lacus purus porta arcu, at ultrices diam nisl eget urna. Curabitur sollicitudin diam quis sollicitudin varius. Ut porta erat ornare laoreet euismod. In tincidunt purus dui, nec egestas dui convallis non. In vestibulum ipsum in dictum scelerisque.

% % Warning text, with a custom title
% \begin{warn}[Notice:]
%   In congue risus leo, in gravida enim viverra id. Donec eros mauris, bibendum vel dui at, tempor commodo augue. In vel lobortis lacus. Nam ornare ullamcorper mauris vel molestie. Maecenas vehicula ornare turpis, vitae fringilla orci consectetur vel. Nam pulvinar justo nec neque egestas tristique. Donec ac dolor at libero congue varius sed vitae lectus. Donec et tristique nulla, sit amet scelerisque orci. Maecenas a vestibulum lectus, vitae gravida nulla. Proin eget volutpat orci. Morbi eu aliquet turpis. Vivamus molestie urna quis tempor tristique. Proin hendrerit sem nec tempor sollicitudin.
% \end{warn}

%----------------------------------------------------------------------------------------

\end{document}
